%% Weak-Memory Local Reasoning, Revisited.
%% Build: latexmk -pdf weak-memory-local-reasoning-revised   (or: tectonic weak-memory-local-reasoning-revised.tex)
\documentclass[11pt]{article}
\usepackage[T1]{fontenc}
\usepackage[margin=1in]{geometry}
\usepackage{microtype}
\usepackage{booktabs}
\usepackage{tabularx}
\usepackage{enumitem}
\usepackage{mathpartir}
\usepackage{xcolor}
\usepackage{placeins}
\usepackage[bookmarks=true,colorlinks=true,linkcolor=black,citecolor=black,urlcolor=blue]{hyperref}

%% The dissertation's own notation, unchanged.
%!TEX root = dissertation.tex

\usepackage{amsmath,amssymb,amsthm}
\usepackage{graphicx}
\usepackage{stmaryrd}
\usepackage{nicefrac}

\theoremstyle{plain}
\newtheorem{lemma}{Lemma}
\newtheorem{theorem}{Theorem}
\newtheorem{corollary}{Corollary}
\newtheorem{proposition}{Proposition}
\newtheorem{conjecture}{Conjecture}

\theoremstyle{definition}
\newtheorem{definition}{Definition}

\theoremstyle{remark}
\newtheorem{note}{Note}
\newtheorem*{remark}{Remark}


\newcommand{\bnfbar}{\ensuremath{~|~}}
\newcommand{\bnfdef}{\ensuremath{~::=~}}

%% set notation
\newcommand{\set}[1]{\left\lbrace #1 \right\rbrace}
\newcommand{\setof}[2]{\set{#1 \mid #2}}
\newcommand{\bigset}[1]{\left\lbrace #1 \right\rbrace}
\newcommand{\bigsetof}[2]{\bigset{#1 \mid #2}}

\DeclareMathOperator{\funcompat}{\smile}

\newcommand{\mtab}{\ensuremath{\;\;\;}}

\newcommand{\powerset}[1]{\ensuremath{\mathcal{P}(#1)}}
\newcommand{\zorch}{\ensuremath{\lightning}}

\DeclareMathOperator{\opdom}{\mathrm{dom}}
\newcommand{\dom}[1]{\ensuremath{\opdom(#1)}}

\newcommand{\bvt}{\ensuremath{\mathbf{t}}}
\newcommand{\bvf}{\ensuremath{\mathbf{f}}}
\newcommand{\llen}[1]{\ensuremath{|#1|}}
\newcommand{\llit}[1]{\ensuremath{\left[#1\right]}}
\newcommand{\lnil}[0]{\ensuremath{\varepsilon}}
\newcommand{\arnil}[0]{\ensuremath{\mathcal{E}}}
\newcommand{\elems}[1]{\ensuremath{|\!|#1|\!|}}
\newcommand{\funof}[1]{\ensuremath{\overline{#1}}}
\DeclareMathOperator{\lapp}{+\!\!\!+}
\DeclareMathOperator{\lcons}{:\!:}
\newcommand{\lhead}[1]{\ensuremath{\mathit{hd}(#1)}}
\newcommand{\ltail}[1]{\ensuremath{\mathit{tl}(#1)}}
\newcommand{\nth}[2]{\ensuremath{#2 ! #1}}
\newcommand{\listof}[1]{\ensuremath{#1~\mathsf{list}}}

\newcommand{\relset}[1]{\ensuremath{|#1|}}
\newcommand{\ext}[1]{\ensuremath{\hat{#1}}}
\newcommand{\restrict}[2]{\ensuremath{#1 |_{#2}}}
\newcommand{\proj}[2]{\ensuremath{#1 |_{#2}}}

\newcommand{\reclit}[1]{\ensuremath{\set{#1}}}

\DeclareMathOperator{\tfun}{\rightarrow}
\DeclareMathOperator{\pfun}{\rightharpoonup}
\DeclareMathOperator{\fpfun}{\pfun_{\mathrm{fin}}}

\DeclareMathOperator{\st}{\,:\,}

\newcommand{\ptup}[2]{\ensuremath{#1 \shortleftarrow #2}}
\newcommand{\funup}[2]{\ensuremath{#1\! \left[ #2 \right]}}
\newcommand{\fundel}[2]{\ensuremath{#1 \setminus#2 }}
\newcommand{\recup}[3]{\funup{#1}{\ptup{#2}{#3}}}
\newcommand{\ptapp}[2]{\ensuremath{#1 \lapp #2}}

\DeclareMathOperator{\conj}{\,\wedge\,}
\DeclareMathOperator{\disj}{\,\vee\,}
\DeclareMathOperator{\onlyif}{\,\Rightarrow\,}
%\DeclareMathOperator{\ifandonlyif}{\,\Longleftrightarrow\,}
\renewcommand{\iff}{\ensuremath{\Leftrightarrow}}

\newcommand{\fresh}[1]{\ensuremath{\mathrm{fresh}(#1)}}
\newcommand{\subst}[2]{\ensuremath{\left[#1 / #2\right]}}
\newcommand{\unsubst}[2]{\ensuremath{\left[#1 \backslash #2\right]}}

\newcommand{\eqclass}[1]{\ensuremath{\left[ #1 \right]}}
\DeclareMathOperator{\downop}{\downarrow}
\newcommand{\downo}[2]{\ensuremath{\downop_{#1}\!#2}}
\newcommand{\down}[1]{\ensuremath{\downop #1}}

%% Calculational proofs

\newcommand{\Calc}[1]{\begin{description}
                \item \begin{tabbing}\qquad\=\quad\=\kill
                \> \> #1\end{tabbing}
                \end{description}}
\newcommand{\Sketch}[1]{\begin{description}
                \item \begin{tabbing}\qquad\=\quad\=\kill
                #1\end{tabbing}
                %\> \> #1\end{tabbing}
                \end{description}}                

\newcommand{\aconn}[1]{\{#1\}\\ }
\newcommand{\aconnl}[1]{\{#1\} }
\newcommand{\aconnt}[1]{\{#1\} $\therefore$\\ }
\newcommand{\aconni}[1]{\> \{#1\}\\ }
\newcommand{\acomm}[1]{\> #1 \\}


\newcommand{\conn}[2]{\\*$#1 $\> \{#2\}\\\> \>}
\newcommand{\Conn}[1]{\conn{=}{#1}}

\newcommand{\ident}[1]{\ensuremath{\mathit{Id}_{#1}}}
\DeclareMathOperator{\rcomp}{\circ}
\newcommand{\relexp}[2]{\ensuremath{\stackrel{#2}{#1}}}
\newcommand{\tcl}[1]{\overset{+}{#1}}
\newcommand{\rtcl}[1]{\overset{\ast}{#1}}

\newcommand{\image}[2]{\ensuremath{#1\left[#2\right]}}
\newcommand{\invimage}[2]{\ensuremath{#1^{-1}\left[#2\right]}}

\newcommand{\lift}[1]{\ensuremath{\mathit{lift}_{#1}}}
\renewcommand{\merge}{\ensuremath{\uplus}}

% \newcommand{\override}{\reflectbox{\rotatebox[origin=c]{180}{\ensuremath{\,\oslash\,}}}}
\newcommand{\override}{\ensuremath{\backslash\!\!\backslash}}

\newcommand{\sentails}{\ensuremath{\,\models\,}}
\newcommand{\satisfies}{\ensuremath{\,\models\,}}
\newcommand{\sequiv}{\ensuremath{\,\equiv\,}}

\newcommand{\fixme}[1]{\todo{FIXME: #1}}

\DeclareMathOperator{\eqdef}{\;=_{\mathit{df}}\;}
\DeclareMathOperator{\iffdef}{\;\equiv_{\mathit{df}}\;}

\newcommand{\card}[1]{\left|#1\right|}
\newcommand{\cmpl}[1]{\ensuremath{\overline{#1}}}

\newcommand{\refines}{\ensuremath{\sqsubseteq}}
\newcommand{\range}[1]{\ensuremath{\mathit{range}(#1)}}

\newcommand{\fun}[1]{\ensuremath{\lambda#1\,.\,}}
%!TEX root = dissertation.tex

\usepackage{stackrel}

%% types

\newcommand{\setlocations}{\ensuremath{\mathbb{L}}}
\newcommand{\setvalues}{\ensuremath{\mathbb{V}}}
\newcommand{\setintegers}{\ensuremath{\mathbb{Z}}}
\newcommand{\setnaturals}{\ensuremath{\mathbb{N}}}
\newcommand{\setpositives}{\ensuremath{\mathbb{N}^+}}
\newcommand{\setidentifiers}{\ensuremath{\mathbb{I}}}
\newcommand{\setbooleans}{\ensuremath{\mathbb{B}}}
\newcommand{\setprocessors}{\ensuremath{\mathbb{P}}}

%% abbreviations

\newcommand{\setstacks}{\ensuremath{\mathbf{Stack}}}
\newcommand{\setheaps}{\ensuremath{\mathbf{Heap}}}
\newcommand{\setbuffers}{\ensuremath{\mathbf{Buf}}}
\newcommand{\setshares}{\ensuremath{\mathbf{Share}}}
\newcommand{\setperms}{\ensuremath{\mathbf{Perm}}}
\newcommand{\setastates}{\ensuremath{\mathit{\Sigma}}}
\newcommand{\setstates}{\ensuremath{\mathbf{State}}}
\newcommand{\setmodels}{\ensuremath{\mathbf{Model}}}
\newcommand{\setmemorysystems}{\ensuremath{\mathbf{Mem}}}
\newcommand{\setpstates}{\ensuremath{\Pi}}
\newcommand{\setmpstates}{\ensuremath{\mathbf{MPState}}}
\newcommand{\setresolvers}{\ensuremath{\mathbf{Resolver}}}
\newcommand{\setconfigurations}{\ensuremath{\mathbf{Config}}}
\newcommand{\setmconfigurations}{\ensuremath{\mathbf{MConfig}}}
\newcommand{\setassertions}{\ensuremath{\mathbf{Asrt}}}

%% operators 

\DeclareMathOperator{\op}{\mathit{op}}
\DeclareMathOperator{\opskip}{\mathsf{skip}}
\DeclareMathOperator{\opfence}{\mathsf{fence}}
\DeclareMathOperator{\opassume}{\mathsf{assume}}
\DeclareMathOperator{\opassert}{\mathsf{assert}}
\DeclareMathOperator{\opassign}{:\!=}
\DeclareMathOperator{\opequals}{=}
\DeclareMathOperator{\opseq}{;}
\DeclareMathOperator{\opchoice}{+}
\DeclareMathOperator{\oppar}{|\!|}
\DeclareMathOperator{\oploop}{\ast}
\DeclareMathOperator{\oplock}{\mathsf{lock}}
\DeclareMathOperator{\opunlock}{\mathsf{unlock}}
\DeclareMathOperator{\opcons}{\mathsf{cons}}
\DeclareMathOperator{\opnew}{\mathsf{new}}
\DeclareMathOperator{\opalloc}{\mathsf{alloc}}
\DeclareMathOperator{\opdispose}{\mathsf{free}}
\DeclareMathOperator{\opif}{\mathsf{if\,}}
\DeclareMathOperator{\opthen}{\mathsf{\,then\,}}
\DeclareMathOperator{\opelse}{\mathsf{\,else}}
\DeclareMathOperator{\opwhile}{\mathsf{while\,}}
\DeclareMathOperator{\opdo}{\mathsf{\,do\,}}
\DeclareMathOperator{\oplocal}{\mathsf{local\,}}
\DeclareMathOperator{\opin}{\mathsf{\,in\,}}
\DeclareMathOperator{\opresource}{\mathsf{resource\,}}
\DeclareMathOperator{\opwith}{\mathsf{with\,}}
\DeclareMathOperator{\opwhen}{\mathsf{\,when\,}}
\DeclareMathOperator{\nil}{\emptyset}
\DeclareMathOperator{\compat}{\smile}


%% primitive commands

\newcommand{\cskip}[0]{\ensuremath{\opskip}}
\newcommand{\cfence}[0]{\ensuremath{\opfence}}
\newcommand{\clock}[0]{\ensuremath{\oplock}}
\newcommand{\cunlock}[0]{\ensuremath{\opunlock}}
\newcommand{\cassume}[1]{\ensuremath{\opassume(#1)}}
\newcommand{\cassert}[1]{\ensuremath{\opassert(#1)}}
\newcommand{\cassign}[2]{\ensuremath{#1 \opassign #2}}
\newcommand{\cload}[2]{\ensuremath{#1 \opassign \left[ #2 \right]}}
\newcommand{\cstore}[2]{\ensuremath{\left[ #1 \right] \opassign #2}}
\newcommand{\cnew}[2]{\ensuremath{#1 \opassign \opnew(#2)}}
\newcommand{\calloc}[1]{\ensuremath{#1 \opassign \opalloc}}
\newcommand{\ccons}[3]{\ensuremath{#1 \opassign \opcons(#2, #3)}}
\newcommand{\cfree}[1]{\ensuremath{\opdispose(#1)}}

%% commands

\newcommand{\cseq}[2]{\ensuremath{#1 \opseq #2}}
\newcommand{\cchoice}[2]{\ensuremath{#1 \opchoice #2}}
\newcommand{\cpar}[2]{\ensuremath{#1 \oppar #2}}
\newcommand{\cifthenelse}[3]{\ensuremath{\opif #1 \opthen #2 \opelse #3}}
\newcommand{\cifthen}[2]{\ensuremath{\opif #1 \opthen #2}}
\newcommand{\cwhile}[2]{\ensuremath{\opwhile #1 \opdo #2}}
\newcommand{\cloop}[1]{\ensuremath{#1^{\oploop}}}
\newcommand{\clocal}[3]{\ensuremath{\oplocal #1 \opassign #2 \opin #3}}
\newcommand{\cres}[2]{\ensuremath{\opresource #1 \opin #2}}
\newcommand{\cwith}[2]{\ensuremath{\left\langle #2\right\rangle_{#1}}}
% \newcommand{\cwith}[2]{\ensuremath{\opwith #1 \opdo #2}}
\newcommand{\catomic}[1]{\ensuremath{\left\langle #1 \right\rangle}}
\newcommand{\ccas}[2]{\ensuremath{\mathsf{cas}_{#1}(#2)}}

\newcommand{\exprs}{\ensuremath{\mathbf{Expr}}}
\newcommand{\fracs}{\ensuremath{\mathbf{Frac}}}
\newcommand{\bexprs}{\ensuremath{\mathbf{BExpr}}}
\newcommand{\pcomms}{\ensuremath{\mathbf{PComm}}}
\newcommand{\comms}{\ensuremath{\mathbf{Comm}}}
\newcommand{\asserts}{\ensuremath{\mathbf{Assert}}}
\newcommand{\progs}{\ensuremath{\mathbf{Prog}}}

\newcommand{\dnexpr}[1]{\ensuremath{\mathcal{E}[\![#1]\!]}}
\newcommand{\dnshare}[1]{\ensuremath{\mathcal{S}[\![#1]\!]}}
\newcommand{\dnbexpr}[1]{\ensuremath{\mathcal{B}[\![#1]\!]}}
\newcommand{\dnpcomm}[1]{\ensuremath{\mathcal{P}[\![#1]\!]}}
\newcommand{\dna}[1]{\ensuremath{\mathcal{A}[\![#1]\!]}}
\newcommand{\dnfrac}[1]{\ensuremath{\mathcal{F}[\![#1]\!]}}

\newcommand{\sflush}[1]{\ensuremath{\mathit{flush}(#1)}}
\newcommand{\speek}[2]{\ensuremath{\mathit{peek}_{#1}(#2)}}
\newcommand{\sval}[1]{\ensuremath{\mathit{val}(#1)}}
\newcommand{\spush}[3]{\ensuremath{\mathit{push}_{#1}^{#2}(#3)}}
\newcommand{\alloc}[1]{\ensuremath{\mathit{alloc}(#1)}}
\newcommand{\allocat}[2]{\ensuremath{\mathit{alloc}_{#1}(#2)}}
\newcommand{\owned}[1]{\ensuremath{\mathit{owned}(#1)}}
\newcommand{\full}[1]{\ensuremath{\mathit{full}(#1)}}

\newcommand{\step}{\ensuremath{\rightarrow}}
\newcommand{\rtstep}{\ensuremath{\stackrel{\ast}{\rightarrow}}}
\newcommand{\nostep}{\ensuremath{\nrightarrow}}
\newcommand{\pstep}[1]{\ensuremath{\underset{#1}{\rightarrow}}}
\newcommand{\taustep}{\pstep{\tau}}
\newcommand{\modstep}{\pstep{\widehat{\tau}}}

\newcommand{\fv}[1]{\ensuremath{\mathrm{fv}(#1)}}
\renewcommand{\mod}[1]{\ensuremath{\mathrm{mod}(#1)}}

\newcommand{\bexpt}{\ensuremath{\mathbf{true}}}
\newcommand{\bexpf}{\ensuremath{\mathbf{false}}}


\newcommand{\spec}[4]{\ensuremath{#1\,\vdash\,\set{#2}\;#3\;\set{#4}}}
\newcommand{\truespec}[4]{\ensuremath{#1\,\models\,\set{#2}\;#3\;\set{#4}}}
\newcommand{\proofof}[1]{\ensuremath{\overline{#1}}}
\newcommand{\andif}{\ensuremath{\text{\;\;\;and\;\;\;}}}
\newcommand{\htriple}[3]{\ensuremath{\set{#1}\;#2\;\set{#3}}}
\newcommand{\triple}[3]{\ensuremath{\vdash\,\set{#1}\;#2\;\set{#3}}}
\newcommand{\truetriple}[3]{\ensuremath{\models\,\set{#1}\;#2\;\set{#3}}}

\newcommand{\safe}[2]{\ensuremath{\mathit{safe}_{#1}(#2)}}

\newcommand{\defined}[1]{\ensuremath{\mathsf{def}(#1)}}

\newcommand{\abort}{\ensuremath{\zorch}}

\newcommand{\live}[1]{\ensuremath{\mathit{live}(#1)}}

\newcommand{\femp}{\ensuremath{\mathbf{emp}}}
\newcommand{\fbar}[1]{\ensuremath{\mathbf{bar}_{#1}}}
\newcommand{\flock}[1]{\ensuremath{\mathbf{lock}_{#1}}}
\newcommand{\fexp}[1]{\ensuremath{\mathbf{exp}({#1})}}
\newcommand{\ftop}[1]{\ensuremath{\mathbf{top}_{#1}}}
%\newcommand{\fwrite}[1]{\ensuremath{\!\stackrel{#1}{\leadsto}\!}}
\newcommand{\fwrite}[1]{\ensuremath{\!\leadsto_{#1}\!}}
\newcommand{\fwrote}[1]{\ensuremath{\rightarrowtail_{#1}}}
\newcommand{\fpointsto}{\ensuremath{\!\mapsto\!}}
\newcommand{\fiter}[1]{\ensuremath{#1^{\ast}}}

\newcommand{\gensym}{\bullet}
\newcommand{\sepsym}{\ast}
\newcommand{\seqsym}{\vartriangleleft}
\newcommand{\sseqsym}{\blacktriangleleft}
\newcommand{\hashsym}{\#}

\DeclareMathOperator{\fhash}{\,\#\,}
\DeclareMathOperator{\fsep}{\,\ast\,}
\DeclareMathOperator{\fseq}{\,\vartriangleleft\,}
\DeclareMathOperator{\fsseq}{\,\blacktriangleleft\,}
\DeclareMathOperator{\fseqc}{\,\vartriangleright\,}
\DeclareMathOperator{\fsseqc}{\,\blacktriangleright\,}
\DeclareMathOperator{\fsepi}{\,\relbar\!\!\!\ast\,}
\DeclareMathOperator{\fseqli}{\,\relbar\!\vartriangleleft\,}
\DeclareMathOperator{\fseqri}{\,\vartriangleleft\!\relbar\,}
\DeclareMathOperator{\fsepic}{\,\ast\!\!\!\relbar\,}
\DeclareMathOperator{\fseqlic}{\,\vartriangleright\!\relbar\,}
\DeclareMathOperator{\fseqric}{\,\relbar\!\vartriangleright\,}
\DeclareMathOperator{\fgen}{\bullet}


\DeclareMathOperator{\shash}{\,\widetilde{\#}\,}
\DeclareMathOperator{\ssep}{\,\widetilde{\ast}\,}
\DeclareMathOperator{\sseq}{\,\widetilde{\vartriangleleft}\,}
\DeclareMathOperator{\ssseq}{\,\widetilde{\blacktriangleleft}\,}
\DeclareMathOperator{\sseqc}{\,\widetilde{\vartriangleright}\,}
\DeclareMathOperator{\ssseqc}{\,\widetilde{\blacktriangleright}\,}
\DeclareMathOperator{\ssepi}{\,\widetilde{\relbar\!\!\ast}\,}
\DeclareMathOperator{\sseqli}{\,\widetilde{\relbar\!\vartriangleleft}\,}
\DeclareMathOperator{\sseqri}{\,\widetilde{\vartriangleleft\!\relbar}\,}
\DeclareMathOperator{\ssepic}{\,\widetilde{\ast\!\!\relbar}\,}
\DeclareMathOperator{\sseqlic}{\,\widetilde{\vartriangleright\!\relbar}\,}
\DeclareMathOperator{\sseqric}{\,\widetilde{\relbar\!\vartriangleright}\,}
\DeclareMathOperator{\sgen}{\widetilde{\bullet}}

\DeclareMathOperator{\mhash}{\,\widehat{\#}\,}
\DeclareMathOperator{\msep}{\,\widehat{\ast}\,}
\DeclareMathOperator{\mseq}{\,\widehat{\vartriangleleft}\,}
\DeclareMathOperator{\msseq}{\,\widehat{\blacktriangleleft}\,}
\DeclareMathOperator{\mseqc}{\,\widehat{\vartriangleright}\,}
\DeclareMathOperator{\msseqc}{\,\widehat{\blacktriangleright}\,}
\DeclareMathOperator{\msepi}{\,\widehat{\relbar\!\!\ast}\,}
\DeclareMathOperator{\mseqli}{\,\widehat{\relbar\!\vartriangleleft}\,}
\DeclareMathOperator{\mseqri}{\,\widehat{\vartriangleleft\!\relbar}\,}
\DeclareMathOperator{\msepic}{\,\widehat{\ast\!\!\relbar}\,}
\DeclareMathOperator{\mseqlic}{\,\widehat{\vartriangleright\!\relbar}\,}
\DeclareMathOperator{\mseqric}{\,\widehat{\relbar\!\vartriangleright}\,}
\DeclareMathOperator{\mgen}{\widehat{\bullet}}


\newcommand{\amod}{\ensuremath{\mathcal{M}}}



\newcommand{\fbaru}{\ensuremath{\mathbf{bar}}}
\newcommand{\fwriteu}{\ensuremath{\!\leadsto\!}}

\newcommand{\pred}[1]{\ensuremath{[\![#1]\!]}}
\newcommand{\locked}[1]{\ensuremath{\mathit{locked}(#1)}}
\newcommand{\taurefines}{\ensuremath{\preceq}}
\newcommand{\modrefines}{\ensuremath{\leq}}
\newcommand{\lockcomplete}[1]{\ensuremath{\text{\textit{lock-complete}}(#1)}}
\newcommand{\stable}[1]{\ensuremath{\mathit{stable}(#1)}}
\newcommand{\stcong}[1]{\ensuremath{\sim_{#1}}}

%% Notation added for this report.
%% Notation that is not in the dissertation.

%% Rocq identifiers: typewriter, underscores allowed, may break at them.
\newcommand{\rocq}[1]{\nolinkurl{#1}}

%% A write by processor i that has flushed: (e ~>_i e') <| bar_i.
%% The dissertation's macro file already defines \fwrote{i} as a tailed
%% arrow with a subscript; we give it this meaning.
\newcommand{\fwroteu}{\ensuremath{\rightarrowtail}}

%% Ownership modes of the repaired judgment.
\newcommand{\modeov}{\ensuremath{\mathsf{ov}}}
\newcommand{\modedj}{\ensuremath{\mathsf{dj}}}
\newcommand{\mspec}[5]{\ensuremath{#1\,\vdash^{#2}\,\set{#3}\;#4\;\set{#5}}}
\newcommand{\mtruespec}[5]{\ensuremath{#1\,\models^{#2}\,\set{#3}\;#4\;\set{#5}}}

%% The lock-free judgment of Section~\ref{sec:daring}.
\newcommand{\bspec}[4]{\ensuremath{#1\,\Vdash\,\set{#2}\;#3\;\set{#4}}}

%% Rule names in running text.
\newcommand{\rn}[1]{\textsc{#1}}

%% Status notes while drafting. Remove before circulating.
\newcommand{\pending}[1]{\textcolor{red}{[\emph{pending:} #1]}}

%% The repaired cell: a write that has flushed, by whichever processor
%% wrote it. The dissertation's all-processor points-to keeps \fpointsto.
\DeclareMathOperator{\fcell}{\!\rightarrowtail\!}

%% Kinds of change, for the differences table.
\newcommand{\kerr}{\textsf{error}}
\newcommand{\kchoice}{\textsf{choice}}
\newcommand{\kforced}{\textsf{forced}}


\title{Weak-Memory Local Reasoning, Revisited}
\author{Ian Wehrman}
\date{Draft of \today}

\begin{document}
\maketitle

\begin{abstract}
My 2012 dissertation proposed a program logic for concurrent programs on an
x86-like machine, where each processor buffers its writes. Its central rule
is a temporal frame rule: to prove that a processor reads back its own
buffered write, a proof sets aside the older contents of the same cell and
reasons about the write alone. I left the soundness proof unfinished, and
out of curiosity I have returned to it with a mechanization in Rocq.

The soundness theorem does not hold as stated. Four of the original rules
derive a specification in which a thread executes $\cskip$ and then owns a
cell with a different value: the proof gives the cell to the lock's
invariant, and the temporal frame rule keeps an invented earlier value as a
claim about the present. However, a repaired calculus is sound and adequate
for the machine. It transfers ownership only when the lock is acquired or
released, and frames temporally only across commands that do not release
the lock, which still covers the read. It keeps the assertion language,
the unit law for $\femp$ and the original load and store axioms. The report
also gives a sound proof system for assertions, worked examples, and a
separate calculus for access without the lock that proves the
dissertation's opening example, a fenced simplification of Dekker's
algorithm. It is written for a reader who knew the dissertation but has
forgotten its details.
\end{abstract}

\section{Introduction}
\label{sec:intro}

In my dissertation~\cite{wehrman12} I proposed a program logic for
concurrent programs on an x86-like machine and left its soundness proof
unfinished, suggesting that the remaining proofs were ``mostly
straightforward inductions.'' I have wondered since whether that was true.
Fourteen years later I returned to the question with the Rocq prover.

The logic was built for programs like this one. A thread owns a memory
cell, writes to it, and immediately reads it back. On this machine the
write first enters the processor's store buffer, and a \emph{flush} later
moves it to the shared heap. If the write flushes between the two
instructions, the read finds the new value in the heap. If it has not, the
heap still holds the old value, and the read finds the new one in the
processor's own buffer. A proof of the read should not have to say which
happened.

Separation logic already lets a proof set aside memory that a command does
not touch. If $R \fsep P$ divides memory by address, the \emph{spatial
frame rule} lets a proof of the command use $P$ alone and add the
\emph{frame} $R$ back afterwards. Here, though, the old value and the
buffered write are at the \emph{same} address. The dissertation added a
second way to divide memory, $R \fseq P$, in which $R$ describes earlier
writes and $P$ later ones, with a \emph{temporal frame rule} to go with it.
For the read, $P$ is the processor's latest write to the cell and $R$ is
everything before it. The load axiom applies to $P$ alone, and the rule
carries $R$ across the load unchanged (Section~\ref{sec:forwarding}).
This rule is what the dissertation's design rests on, and it is also what
breaks its soundness.

\paragraph{A proof can give a cell away and keep it.}
The failure needs no buffered writes. Take a thread that holds the lock and
owns a cell containing $7$, and let the lock's \emph{invariant}, the
assertion describing the memory it protects, allow either nothing or that
one cell. The original invariant rule lets a proof of $\cskip$ hand the cell
over to the invariant. Nothing happens on the machine, but the proof no
longer counts the cell as the thread's.

Before the handover, the proof can also set aside an invented history of
the cell. A temporal division is checked against memory as it now stands,
and where both parts contain the same cell, the later part's value is the
one in memory. So the thread's state also satisfies ``the cell contained
$5$, and later the cell contained $7$,'' although $5$ was never written.
The temporal frame rule carries the earlier part across $\cskip$. Afterwards
the later part holds only the lock, since its cell went to the invariant,
and the unit law $R \fseq \femp \sentails R$ turns what is left into a
claim about the present: the thread owns the cell, and it contains $5$. The
heap still holds $7$. Four of the original rules derive this
(Section~\ref{sec:inconsistent}), so no choice of semantics can make them
sound. Several auxiliary claims of the dissertation are also false, but each
can be fixed without changing the design (Section~\ref{sec:local-errors}).

\paragraph{The cause is general.}
Nothing in this argument depends on store buffers. An unrestricted temporal
frame rule, a right unit for temporal conjunction, the rule of consequence,
and a command that gives away a cell cannot coexist once some history that
held before the transfer is false after it (Section~\ref{sec:trilemma}). A
sequential conjunction in which later values override earlier ones supplies
such histories.

\paragraph{A repaired calculus is sound.}
In the repair, a thread receives the invariant's cells only when it acquires
the lock, and gives cells up only when it releases the lock, saying which
ones it returns to satisfy the invariant. That removes the handover during
$\cskip$. A release still gives cells away, so the temporal frame rule
applies only to commands that contain no $\cunlock$. A load is such a
command, and the read above is still proved locally. Four side conditions
change as well (Section~\ref{sec:changes}). With two distinct active
processors and an invariant from a restricted fragment, the repaired
calculus is sound, and it is adequate for the dissertation's machine: a
verified program started as its specification requires cannot abort, and if
it terminates, its postcondition holds (Section~\ref{sec:results}). The
proofs are checked by the Rocq kernel and use no axioms.

\paragraph{Most of the design survives.}
The machine, the assertion language and its interpretation of atoms, the
unit laws for $\femp$, overlapping temporal history, the temporal and
spatial frame rules, the concurrency rule, and the original load and store
axioms are all retained. So is most of the dissertation's assertion
algebra. The exceptions involve its third conjunction, the strong temporal
conjunction $P \fsseq Q$, which asks for a spatial and a temporal division
of memory that need not agree. The laws that rebracket or exchange through
it fail, and with them its frame rule. The exchange laws between spatial
and temporal conjunction also fail or are vacuous
(Appendix~\ref{app:statements}).

\paragraph{Assertions have a proof system.}
A calculus of 47 rules for entailment between assertions is proved sound
(Section~\ref{sec:assertion-rules}). Its completeness was not a goal, but it
is enough: the larger derivations below use no semantic entailment at all.

\paragraph{The logic can be used.}
Section~\ref{sec:proofs} derives a shared container in which a thread
acquires the lock, removes a cell, releases the lock, writes the cell
privately, and later returns it, while another processor's write stays
buffered throughout. It re-derives the dissertation's own message-passing
example. And it builds named locks on top of the single global lock, as
derived rules, and uses one to pass a message of any length from a producer
to a consumer.

\paragraph{Unlocked access has a sound rule.}
The dissertation's rules for touching shared state without the lock are
unsound for the repaired calculus. Section~\ref{sec:daring} gives a
restricted form, in a separate judgment for programs that take no lock, and
uses it to prove the fenced simplification of Dekker's algorithm with
which the dissertation opens.

\medskip
This report is for a reader who knew the original dissertation but has
forgotten its details. Section~\ref{sec:original} recalls the original model
and calculus in the original notation. Sections~\ref{sec:wrong}
and~\ref{sec:changes} say what was wrong and what changed.
Sections~\ref{sec:model} to~\ref{sec:proofs} present the repaired model,
calculus, results and derivations. Proofs are not reproduced. The text
names the principal Rocq theorems, and the appendix gives fuller
statements.

\section{The original model and calculus}
\label{sec:original}

This section recalls Chapter~4 of the dissertation, in its notation. Nothing
here is new. Page and figure numbers refer to the dissertation.

\subsection{The machine}

A \emph{memory system} is a triple $\mu = (h,B,K)$: a heap
$h : \setlocations \pfun \setvalues$, an array of write buffers
$B : \setprocessors \tfun \listof{(\setlocations \times \setvalues)}$ with
the oldest write first, and a set $K \subseteq \setprocessors$ of
\emph{blocked} processors. A state pairs a stack $s$ with a memory system.
The lock is not a separate component: it is free when $K = \nil$ and held
by $i$ when $K = \setprocessors \setminus \set{i}$.

Primitive commands carry the identifier of the processor that executes
them. Their behaviour on processor $i$ is as follows.
\begin{itemize}[nosep]
  \item $\cstore{e}{e'}_i$ appends a write to $B(i)$. The address must be
        allocated, in the heap or in $B(i)$. A blocked processor may still
        store.
  \item $\cload{x}{e}_i$ requires $i \notin K$ and returns the newest write
        to that address in $B(i)$, or else the heap value.
  \item $\cfence_i$, $\clock_i$ and $\cunlock_i$ require $B(i) = \lnil$.
        $\clock_i$ additionally requires $K = \nil$ and sets
        $K = \setprocessors \setminus \set{i}$; $\cunlock_i$ reverses this.
  \item A silent step $\taustep$ moves the oldest write of some unblocked
        processor from its buffer to the heap.
\end{itemize}
Commands are built from primitives by sequencing $\cseq{c}{c'}$, choice
$\cchoice{c}{c'}$, iteration $\cloop{c}$ and parallel composition
$\cpar{c}{c'}$. There is no allocation. While a processor holds the lock,
no other processor can load or flush, but others may continue to store and
to assign to registers.

\subsection{Generalized memories and separation}

A \emph{generalized memory system} $\nu = (h,B,K,\Gamma)$ adds a
\emph{buffer-completeness} set $\Gamma \subseteq \setprocessors$, with
$h = \nil$ whenever $\Gamma = \nil$. Membership $i \in \Gamma$ means that
buffer $i$ is complete: nothing may be placed before it.

Three operations combine generalized memories. Write $\nu_k =
(h_k,B_k,K_k,\Gamma_k)$. All three take $K_0 \cup K_1$ and
$\Gamma_0 \cup \Gamma_1$ and require $K_0 \cap K_1 = \nil$.
\begin{description}[nosep,leftmargin=1.5em]
  \item[Spatial, $\nu_0 \msep \nu_1$.] The address domains of the two
    sides, counting both heaps and buffers, are disjoint. The heap is
    $h_0 \uplus h_1$ and each buffer is an interleaving of the two.
  \item[Temporal, $\nu_0 \mseq \nu_1$.] The heap is $h_0 \override h_1$, so
    later values win, and each buffer is $B_0(i) \lapp B_1(i)$. The sides
    may share addresses, subject to a compatibility condition: an unblocked
    pending write in $\nu_0$ may not target an address in the heap of
    $\nu_1$, nor an address pending on a different unblocked processor in
    $\nu_1$. Finally $B_0(i) = \lnil$ for every $i \in \Gamma_1$.
  \item[Spatiotemporal, $\nu_0 \mhash \nu_1$.] As temporal, but each buffer
    is any interleaving whose lookup function agrees with
    $B_0(i) \lapp B_1(i)$, and the condition on $\Gamma_1$ is dropped. It
    contains both of the others and is used only in the semantics of
    specifications.
\end{description}
\emph{Strong temporal} separation $\nu_0 \msseq \nu_1$ is $\nu_0 \mseq
\nu_1$ restricted to pairs that are also spatially compatible.

\subsection{Assertions}

Assertions are given by
\begin{align*}
  P \bnfdef{} & b \bnfbar P \disj P \bnfbar P \conj P \bnfbar \exists x \st P
    \bnfbar \forall x \st P \bnfbar \femp \bnfbar \fbar{e} \bnfbar
    \flock{e} \bnfbar e \fwrite{e'} e'' \\
  \bnfbar{} & P \fsep P \bnfbar P \fseq P \bnfbar P \fsepi P \bnfbar
    P \fseqli P \bnfbar P \fseqri P
\end{align*}
and are interpreted over a stack and a generalized memory. The atoms are:
\begin{itemize}[nosep]
  \item $\femp$: empty heap, empty buffers, $K = \nil$; $\Gamma$ is
        unconstrained.
  \item $\fbar{i}$: as $\femp$, with $i \in \Gamma$ when $i$ is a processor.
  \item $\flock{i}$: empty heap and buffers, with
        $K = \setprocessors \setminus \set{i}$.
  \item $e \fwrite{i} e'$ (``leads to''): either a single \emph{pending}
        write $(e,e')$ in $B(i)$ with an empty heap and
        $K \subseteq \set{i}$, or a single \emph{flushed} heap cell
        $e \mapsto e'$ with empty buffers, $K = \nil$ and $i \in \Gamma$.
\end{itemize}
$P \fsep Q$ and $P \fseq Q$ hold of $\nu$ when $\nu$ is a spatial,
respectively temporal, combination of a model of $P$ and a model of $Q$.
The three implications are the corresponding adjoints. Two abbreviations
matter below:
\[
  P \fsseq Q \eqdef (P \fsep Q) \conj (P \fseq Q)
  \qquad\qquad
  e \fpointsto e' \eqdef (\exists x \st e \fwrite{x} e') \fseq
                         (\forall x \st \fbar{x}).
\]
Points-to therefore describes a flushed cell with \emph{every} processor
buffer-complete. The dissertation claims that every assertion is closed
under flushing (Proposition~9), and that each conjunction is associative
with $\femp$ as a unit (Figure~4.5).

\subsection{Specifications and rules}

A specification $\spec{J}{P}{c}{Q}$ has an invariant $J$ for shared state
and a private precondition and postcondition. Figure~\ref{fig:printed}
reproduces the axioms and rules.

\begin{figure}[t]
\small
\begin{mathpar}
\inferrule[skip]{ }{\spec{J}{P}{\cskip_i}{P}}
\and
\inferrule[assume]{ }{\spec{J}{\neg b \disj P}{\cassume{b}_i}{P}}
\and
\inferrule[assert]{ }{\spec{J}{b \conj P}{\cassert{b}_i}{P}}
\and
\inferrule[assign]{ }{\spec{J}{P\subst{e}{x}}{\cassign{x}{e}_i}{P}}
\and
\inferrule[fence]{ }{\spec{J}{\femp}{\cfence_i}{\fbar{i}}}
\and
\inferrule[lock]{ }{\spec{J}{\femp}{\clock_i}{\flock{i}}}
\and
\inferrule[unlock]{ }{\spec{J}{\flock{i}}{\cunlock_i}{\femp}}
\and
\inferrule[load]{ }
  {\spec{J}{e \fwrite{i} e' \fsseq P}{\cload{x}{e}_i}
        {(e \fwrite{i} e' \fsseq P) \conj x = e'}}
\and
\inferrule[store]{ }
  {\spec{J}{e \fwrite{i} e'' \fsseq P}{\cstore{e}{e'}_i}
        {(e \fwrite{i} e'' \fsseq P) \fseq e \fwrite{i} e'}}
\and
\inferrule[frame-$\bullet$]
  {\spec{J}{P}{c}{Q} \\ \mod{c} \cap \fv{R} = \nil}
  {\spec{J}{R \fgen P}{c}{R \fgen Q}}
\and
\inferrule[cons]
  {P \sentails P' \\ \spec{J}{P'}{c}{Q'} \\ Q' \sentails Q}
  {\spec{J}{P}{c}{Q}}
\and
\inferrule[seq]
  {\spec{J}{P}{c}{R} \\ \spec{J}{R}{c'}{Q}}
  {\spec{J}{P}{\cseq{c}{c'}}{Q}}
\and
\inferrule[choice]
  {\spec{J}{P}{c}{Q} \\ \spec{J}{P}{c'}{Q}}
  {\spec{J}{P}{\cchoice{c}{c'}}{Q}}
\and
\inferrule[loop]{\spec{J}{P}{c}{P}}{\spec{J}{P}{\cloop{c}}{P}}
\and
\inferrule[disj]
  {\spec{J}{P}{c}{Q} \\ \spec{J}{P'}{c}{Q}}
  {\spec{J}{P \disj P'}{c}{Q}}
\and
\inferrule[ex]
  {\spec{J}{P}{c}{Q} \\ x \notin \fv{c,Q}}
  {\spec{J}{\exists x \st P}{c}{Q}}
\and
\inferrule[conc]
  {\spec{J}{P_k}{c_k}{Q_k} \\ \fv{P_k,c_k,Q_k} \cap \mod{c_{1-k}} = \nil
   \quad (k = 0,1)}
  {\spec{J}{P_0 \fsep P_1}{\cpar{c_0}{c_1}}{Q_0 \fsep Q_1}}
\and
\inferrule[share]
  {\spec{J}{P}{c}{Q}}
  {\spec{\femp}{J \fsep P}{c}{J \fsep Q}}
\and
\inferrule[inv]
  {\spec{\femp}{J \fsep P \fsep \flock{i}}{c}{J \fsep Q \fsep \flock{i}}}
  {\spec{J}{P \fsep \flock{i}}{c}{Q \fsep \flock{i}}}
\end{mathpar}
\caption{The printed axioms and rules (Figures~4.7 and~4.8).
\rn{frame-$\bullet$} stands for three rules, \rn{frame-sp}, \rn{frame-tm}
and \rn{frame-stm}, with $\bullet$ read as $\fsep$, $\fseq$ and $\fsseq$.}
\label{fig:printed}
\end{figure}

Three features of the design are worth recalling, because the repair turns
on them. First, \rn{lock} and \rn{unlock} produce and consume only the
token $\flock{i}$. The invariant reaches the lock holder through \rn{inv},
which may be applied to any command executed while the lock is held and
allows the private state to exchange cells with $J$. Second, \rn{frame-tm}
has only the usual variable condition: any history $R$ may be placed before
any specification. Third, temporal conjunction is not cancellative, as the
dissertation remarks: $\mu \sseq \mu_0 = \mu \sseq \mu_1$ does not imply
$\mu_0 = \mu_1$.

\subsection{Semantics and status}

Truth of a specification is defined through a family of predicates
$\safe{n}{c,s,\mu,J,Q}$ after Vafeiadis. The predicate holds at $n+1$ when:
if $c$ is $\cskip$ then $(s,\mu,\setprocessors) \sentails Q$; and for every
$\mu_J$, $\mu_F$ and $\mu_0 \in \mu_J \ssep (\mu_F \shash \mu)$ that is
lock-complete, with either $(s,\mu_J,\setprocessors) \sentails J$ or the
lock held in $\mu_0$, the configuration does not abort, and every step to
$(c',s',\mu_1)$ admits a new decomposition
$\mu_1 \in \mu'_J \ssep (\mu'_F \shash \mu')$ with
$\mu'_F \taurefines \mu_F$, the same condition on $J$, and
$\safe{n}{c',s',\mu',J,Q}$. A specification is true,
$\truespec{J}{P}{c}{Q}$, when $(s,\mu,\setprocessors) \sentails P$ implies
$\safe{n}{c,s,\mu,J,Q}$ for all $n$. Local states are always read at full
completeness, $\Gamma = \setprocessors$.

Theorem~1 states that provable specifications are true. The dissertation
gives soundness lemmas for the axioms in an appendix, sketches a frame lemma for the
spatiotemporal operator $\shash$, and reduces the spatial frame rule to
Conjecture~1, which is left open. Associativity of buffer overriding
(Proposition~2) and of the three separations, and flushing closure of
assertions (Proposition~9), are stated with at most a sketch.

\section{What was wrong}
\label{sec:wrong}

The problems are presented from the smallest to the most serious,
which is also the order in which the mechanization found them. Each is a
theorem checked by Rocq, and the theorem's name is given so that the
statement can be inspected. All of them are stated against a transcription of the original definitions.

\subsection{Errors that can be fixed in place}
\label{sec:local-errors}

Table~\ref{tab:local} lists six original claims that are false as stated.
None is deep: each is fixed by correcting the one definition or side
condition involved, and the rest of the design is unaffected.



\begin{table}[!ht]
\small
\begin{tabularx}{\textwidth}{@{}p{0.27\textwidth}X p{0.23\textwidth}@{}}
\toprule
Original claim & Counterexample & Repair \\
\midrule
\rn{assign} is sound for every invariant
  & With $J$ the assertion $x = 0$, the axiom derives
    $\spec{J}{\bexpt}{\cassign{x}{1}_1}{\bexpt}$, and the assignment
    falsifies $J$ with the lock free.
  & Require $x \notin \fv{J}$. \\[2pt]
Buffer overriding is associative (Prop.~2)
  & Three one-element buffers writing $1$, $2$, $3$ to one address: the
    list $[2,1,3]$ is in $a \override (b \override c)$ but not in
    $(a \override b) \override c$.
  & A write from the later buffer may be taken only after the earlier
    buffer's writes to that address. \\[2pt]
Spatiotemporal separation is associative
  & Compatibility ignores blocked processors, so whether two memories
    compose can depend on a third that supplies the blocking fact.
  & Check conflicts whether or not the processor is blocked. \\[2pt]
Assertions are closed under flushing (Prop.~9)
  & $\bexpf \fseqri \fbar{1}$ holds vacuously while processor $1$ has a
    pending write, and fails once the write flushes.
  & Implications quantify over all futures of the current memory. \\[2pt]
Conjecture~1: after a step, a spatiotemporal split of a spatially
split memory is spatial, if its first part is a flushed descendant of the
original first part
  & A cell duplicated across $\shash$ with equal values is not a spatial
    split.
  & Not needed by the final proof. \\[2pt]
\rn{unlock} is sound for every invariant
  & With $J = \bexpf$ and the lock held, $\cunlock_i$ is enabled and
    nothing restores $J$.
  & See Section~\ref{sec:changes}. \\
\bottomrule
\end{tabularx}
\caption{Original claims that fail, each with a checked counterexample.
Appendix~\ref{app:definitions} states the three repaired definitions.}
\label{tab:local}
\end{table}

Two remarks. The repaired compatibility relation is stronger than the
original one, so it rejects some compositions the dissertation admits. Those
are states in which a blocked processor has a pending write to a cell that
someone else owns. And the repaired implications are
the ordinary Kripke reading of implication over the flushing preorder; the
three adjunctions between conjunctions and implications still hold.

The failure of \rn{unlock} points at something larger. The original safety
predicate asks for the invariant only while the lock is free. Once a
thread holds the lock, the predicate has forgotten that $J$ held at the
moment of acquisition, so neither \rn{unlock} nor \rn{inv} can be justified
from it; \rn{inv} on its own takes a true specification to a false one. That alone would call for a different predicate. The next
subsection shows that it calls for different rules.

\subsection{The original rules are inconsistent}
\label{sec:inconsistent}

The derivations below use the dissertation's own points-to, $x \fpointsto
v$, and are checked against its own interpretation of assertions.

\paragraph{History can be invented.}
A model is a memory and records nothing of how that memory came about. So a
current cell is also a model of any stale value followed by that cell:
\begin{equation}
  \label{eq:invent}
  x \fpointsto 7 \fsep \flock{i}
  \;\sentails\;
  x \fpointsto 5 \fseq (x \fpointsto 7 \fsep \flock{i}).
\end{equation}
The stale cell is overridden by the later one, and nothing remains of it.
This is the non-cancellativity that the dissertation notes in passing.

\paragraph{A $\cskip$ that changes memory.}
Let $J$ be $\femp \disj x \fpointsto 7$. The following is a derivation in
the original calculus.
\begin{enumerate}[nosep]
  \item $\spec{\femp}{J \fsep x \fpointsto 7 \fsep \flock{i}}{\cskip_i}
        {J \fsep \femp \fsep \flock{i}}$,
        by \rn{skip} and \rn{cons}: $J$ cannot be its second disjunct, since
        the cell would occur twice, and in the first disjunct the cell can
        be reassigned to $J$.
  \item $\spec{J}{x \fpointsto 7 \fsep \flock{i}}{\cskip_i}
        {\femp \fsep \flock{i}}$, by \rn{inv}. The thread has given its
        cell to the invariant.
  \item $\spec{J}{x \fpointsto 5 \fseq (x \fpointsto 7 \fsep \flock{i})}
        {\cskip_i}{x \fpointsto 5 \fseq (\femp \fsep \flock{i})}$,
        by \rn{frame-tm}.
  \item $\spec{J}{x \fpointsto 7 \fsep \flock{i}}{\cskip_i}
        {x \fpointsto 5 \fsep \flock{i}}$, by \rn{cons}, using
        \eqref{eq:invent} on the left and the unit law for $\femp$ on the
        right.
\end{enumerate}
The conclusion says that a thread holding the lock and a cell containing
$7$ may do nothing and then own that cell containing $5$. In the state
where processor $1$ holds the lock and $x$ contains $7$, the precondition
holds and the postcondition does not.
(\rocq{original_points_stale_value_unsound}; the mechanized statement parses
\rn{inv}'s premise as $J \fsep (P \fsep \flock{i})$.)

\paragraph{A $\cskip$ that drains another processor's buffer.}
Spatial conjunction takes the union of the completeness sets and does not
require them to be disjoint, so an empty neighbour can copy a cell's marks:
$x \fpointsto 7 \fsep \femp \sentails x \fpointsto 7 \fsep \fbar{2}$. With
$J$ the cell $x \fpointsto 7$, \rn{skip}, \rn{cons} and \rn{inv} give
$\spec{J}{\femp \fsep \flock{1}}{\cskip_1}{\fbar{2} \fsep \flock{1}}$, and
\rn{frame-tm} with a pending write $R = y \fwrite{2} 17$ gives
\[
  \spec{J}{R \fseq (\femp \fsep \flock{1})}{\cskip_1}
          {R \fseq (\fbar{2} \fsep \flock{1})}.
\]
When processor $1$ holds the lock and processor $2$'s write to $y$ is still
buffered, the precondition holds. The postcondition places a barrier for
processor $2$ after its pending write, which that memory does not satisfy.
(\rocq{original_points_foreign_barrier_unsound}.)

\paragraph{The strong temporal frame rule does not escape.}
One might hope that \rn{frame-stm} is safe, since its frame is also
spatially separate. It is not, because $R \fsseq P$ abbreviates
$(R \fsep P) \conj (R \fseq P)$ and the two conjuncts may be witnessed by
different splits. Take $R = \femp \disj x \fpointsto 5$. In the state above,
the spatial conjunct of $R \fsseq (x \fpointsto 7 \fsep \flock{1})$ holds
with $R$ empty, and the temporal conjunct holds with $R$ the stale cell.
Applying \rn{frame-stm} to step~2 yields a postcondition that the unchanged
state cannot satisfy. (\rocq{original_strong_frame_counterexample}; this
derivation uses \rn{inv} and \rn{frame-stm} but not \rn{frame-tm}.)

All three conclusions are also invalid for the original safety predicate,
but that is the lesser point. They are false under the
dissertation's informal reading of a specification, and under any semantics
for which a true specification describes what the program does.

\subsection{Why no semantics could repair it}
\label{sec:trilemma}

The first derivation is an instance of a general conflict.

\begin{proposition}[\rocq{trilemma}]
\label{prop:trilemma}
Suppose a logic of triples has an unrestricted temporal frame rule, the
rule of consequence, and the entailment $R \fseq \femp \sentails R$.
Suppose $\htriple{P}{c}{\femp}$ is provable. Then $\htriple{R \fseq P}{c}{R}$
is provable for every $R$. Hence if some execution of $c$ starts in a state
satisfying $R \fseq P$ and ends in a state that does not satisfy $R$, the
logic is unsound.
\end{proposition}

The proof is two rule applications. The content is in the last sentence. A
triple $\htriple{P}{c}{\femp}$ with $P$ owning a cell is a \emph{transfer-out
contract}: after $c$, the cell belongs to someone else. Whenever temporal
conjunction lets later values override earlier ones, $R$ may describe an
older state of the same cell, as the store axiom's own postcondition does.
Framing $R$ across the transfer then hands the thread a stale copy of a
cell it no longer owns.

So four things cannot all be kept: the unrestricted temporal frame rule,
the right-unit law, transfer-out contracts ending in $\femp$, and
overlapping temporal history. A sound logic must give one up, and the
choice determines its character:
\begin{itemize}[nosep]
  \item \emph{No transfer out.} Cells may be taken from the invariant and
        never returned. This is sound and useless for a shared container.
  \item \emph{No right unit.} Let $\femp$ record that a cell was given
        away. Framing survives, and every postcondition that mentions
        $\femp$ changes meaning.
  \item \emph{Transfer leaves a mark.} Return an explicit tombstone
        assertion instead of $\femp$. The algebra survives; every later
        specification carries the tombstones.
  \item \emph{Restrict the frame.} Forbid temporal frames that overlap the
        cells being given away. Everything else is kept.
\end{itemize}
The first two were built and abandoned during this work, for the reasons
given. The repaired calculus takes the fourth.
Proposition~\ref{prop:trilemma} shows that some restriction of this kind is
needed, and the derivations above supply histories that do become false
after a transfer. It does not show that the particular restriction adopted
is the weakest possible.

The second derivation has a different cause and a different repair. The
all-processor completeness of points-to, and more generally the
completeness marks of whatever the invariant holds, should not pass to a
thread that acquires the lock. A thread can vouch only for its own buffer.

\subsection{The dissertation's own example}
\label{sec:own-example}

Section~4.1 of the dissertation proves a small message-passing program
correct, using an invariant that mentions two registers and the derived
rule \rn{atomic}. That proof does not carry over to the repaired calculus:
\rn{atomic} and \rn{share} in their original forms, and
auxiliary-variable elimination, are not available, and the concurrency rule now forbids either thread from
modifying a variable free in the invariant. The proof does not carry
over, although its conclusion is still true.
Section~\ref{sec:own-example-repaired} derives the same program from a
different invariant and recovers the original conclusion.

\section{What changed}
\label{sec:changes}

Table~\ref{tab:changes} lists every difference between the original calculus
and the one proved sound. Each is marked as an \kerr{} (the original
statement is false, with a checked counterexample), as \kforced{} by an
error elsewhere, or as a \kchoice{} among sound alternatives.

\begin{table}[!ht]
\small
\begin{tabularx}{\textwidth}{@{}p{0.33\textwidth}Xl@{}}
\toprule
Original & Repaired & Kind \\
\midrule
\rn{assign}, \rn{load}, \rn{ex}, \rn{conc} say nothing of $\fv{J}$
  & The assigned, loaded or bound variable, and every variable a parallel
    branch modifies, is absent from $\fv{J}$ & \kerr \\
List overriding by lookup; compatibility ignores blocked processors;
pointwise implications
  & Ordered overriding; stable compatibility; implications quantify over
    futures (Table~\ref{tab:local}) & \kerr \\
\rn{lock} and \rn{unlock} exchange a token; \rn{inv} moves cells
  & \rn{inv} is removed; \rn{lock} hands over $J$ and \rn{unlock} takes it
    back & \kerr \\
\rn{frame-tm} for every specification
  & \rn{frame-tm} only for commands that contain no $\cunlock$ & \kerr \\
\rn{frame-stm}
  & Omitted. Sound for one common split of the frame, unsound for
    $(R \fsep P) \conj (R \fseq P)$ & \kerr \\
\rn{share}
  & Holds as a theorem outside the calculus, and only when the
    precondition and postcondition are lock-free & \kforced \\
$e \fpointsto e'$ with every processor complete
  & $e \fcell e'$, complete for some processor; load and store rules for
    it & \kforced \\
Any assertion as invariant
  & Invariants from a guarded, buffer-free fragment & \kchoice \\
Any set of processors
  & At least two active processors & \kforced \\
\bottomrule
\end{tabularx}
\caption{Differences between the original and the repaired calculus.}
\label{tab:changes}
\end{table}

Three of these need a word of explanation.

\paragraph{Transferring lock rules.}
Section~\ref{sec:local-errors} observed that the original predicate forgets
the invariant at acquisition. The standard cure, from concurrent separation
logic, is to move ownership with the lock: acquiring gives the thread the
cells of $J$, and releasing requires it to give back cells satisfying $J$.
\rn{inv} then has nothing left to do. Removing it is also what blocks the
three derivations of Section~\ref{sec:inconsistent}, all of which are
proofs about $\cskip$: restricting the frame rules alone would not have
been enough.

\paragraph{The cell.}
A thread that acquires the lock can vouch for its own buffer, which the
machine required to be empty, and for no one else's. So the cells it
receives carry its completeness mark only. The
original points-to demands every processor's mark, which one acquisition
cannot supply, so a general \rn{lock} rule cannot promise it. The repaired calculus uses
\[
  e \fcell e' \eqdef \exists x \st (e \fwrite{x} e' \fseq \fbar{x}),
\]
a flushed cell that is complete for whichever processor is credited with
the write. This is the definition the dissertation considered first and
rejected, on the grounds that $1 \fwrite{0} 1 \fseq 2 \fcell 2$ would then
have models in which the write to location $1$ is still buffered. Those
models are reachable, stable states of the machine. Because a model does
not record who wrote a flushed cell, any marked processor may be credited,
and two further axioms let any processor load from or store to such a
cell. The original points-to remains expressible and keeps its meaning.

\paragraph{Two processors.}
The lock is represented only by the set of blocked processors. With a
single processor there is no one to block: a held lock is indistinguishable
from a free one, $\flock{i}$ is equivalent to $\femp$, and
$\cseq{\clock_i}{\clock_i}$ runs to completion. A rule that hands over the
invariant at acquisition would then hand it over twice. On one processor the calculus derives a false
specification of $\cseq{\clock_i}{\clock_i}$. The assumption is
not about finiteness, and the second processor need not execute anything.

\paragraph{Roads not taken.}
Two other repairs were built during this work and abandoned, and they are
the first two exits of Section~\ref{sec:trilemma}. Keeping the invariant's
assertions in a separate compartment of the judgment recovers the
unrestricted \rn{frame-tm}, but a cell can then never be returned to the
invariant. Letting $\femp$ record cells that were given away permits the
return, but $R \fseq \femp$ no longer entails $R$, and in that
interpretation an ordinary later cell can silently erase an earlier,
unrelated one.

\section{The repaired model}
\label{sec:model}

The machine is unchanged. What changes is the predicate that says when a
command is safe.

\paragraph{Resources.}
A thread's local state is a generalized memory $\nu$ with its \emph{actual}
completeness set. The original predicate read every local state at
$\Gamma = \setprocessors$, which is why it could not distinguish a barrier
the thread had earned from one it had not. A mark $i \in \Gamma$ now means
what the dissertation said informally: nothing of processor $i$'s is
pending before this resource. A program step may add a mark
only for its own processor, and only if it is $\cfence_i$, $\clock_i$ or
$\cunlock_i$, whose guards require the whole of buffer $i$ to be empty.
Separately, the futures of a resource, under which safety is closed,
include flushes of its own writes and growth of its marks; those concern
temporal composition within the resource and do not certify that a buffer
of the machine has drained.

\paragraph{Environments.}
The rest of the machine is described by a \emph{pool} $\mu_J$, holding the
invariant's cells while the lock is free, and an external \emph{frame}
$\nu_F$, standing for other threads and for any history that a frame rule
has set aside. An environment for $\nu$ in mode $m$ is a pool, a frame and
a machine memory $\mu_0$ with
\[
  \mu_0 \in \mu_J \ssep (\nu_F \mhash \nu)
\]
such that either the lock is free in $\mu_0$ and
$(s,\mu_J,\setprocessors) \sentails J$, or the lock is held and $\mu_J$ is
empty. In the held case the token may be in $\nu$, in $\nu_F$ when another
thread holds the lock, or split between them. The mode restricts the frame:
\begin{itemize}[nosep]
  \item in \emph{overlap mode}, $\modeov$, any compatible frame is allowed,
        including one that mentions addresses of $\nu$;
  \item in \emph{disjoint mode}, $\modedj$, the frame and $\nu$ share no
        address.
\end{itemize}

\paragraph{Safety.}
$\safe{n+1}{c,s,\nu,J,Q}$ holds in mode $m$ when, for every future
$\nu_1 \modrefines \nu$:
\begin{enumerate}[nosep]
  \item if $c$ is $\cskip$ then $s,\nu_1 \sentails Q$; and
  \item for every environment $(\mu_J,\nu_F,\mu_0)$ of $\nu_1$, the
        configuration $c,(s,\mu_0)$ does not abort, and for every
        \emph{program} step to $c',(s',\mu_1)$ there are $\nu'$ and
        $\mu'_J$ forming an environment of $\nu'$ with the \emph{same}
        frame $\nu_F$ and memory $\mu_1$, such that $\safe{n}{c',s',\nu',J,Q}$
        and the step's bookkeeping conditions hold.
\end{enumerate}
The bookkeeping conditions say that the pool changes only at $\clock$ and
$\cunlock$; that $\cunlock_i$ requires the whole token in $\nu_1$; that the
addresses of $\mu'_J$ and $\nu'$ lie within those of $\mu_J$ and $\nu_1$;
and that $\nu'$ gains no completeness mark beyond the one the step earns.
A specification is valid in mode $m$, written $\mtruespec{J}{m}{P}{c}{Q}$,
when $s,\nu \sentails P$ implies safety at every $n$. Overlap-mode validity
implies disjoint-mode validity. Appendix~\ref{app:safety} gives these
definitions in full.

Four features distinguish this from the original predicate.
\begin{itemize}[nosep]
  \item \emph{The frame is fixed within a step and arbitrary between
        steps.} The original predicate let the frame take silent steps of
        its own choosing during a program step. That permits a local
        residual to claim a pending write that belonged to the frame, and
        the frame lemma fails.
        Re-quantifying the frame at every step is what lets another
        thread, including one holding the lock, act in between.
  \item \emph{Flushes are handled by closure.} A silent step is attributed
        to the frame or to the local resource with the other fixed, and safety is closed
        under futures of the local resource by its outermost quantifier.
  \item \emph{The mark budget.} Without it, a proof could choose a
        successor that claims another processor's barrier.
  \item \emph{The modes.} \rn{unlock} is valid only in disjoint mode: the
        released cells become the pool, and a frame holding a stale copy of
        one would collide with it. That is Proposition~\ref{prop:trilemma}
        seen from the semantics.
\end{itemize}

\paragraph{Three lemmas carry the proofs.}
With the frame fixed, the local residual's buffers are determined by the
machine's: ordered overriding places the frame's writes to an address
before the local ones, so there is one way to subtract the frame. Every machine flush belongs to exactly
the frame or the local resource. And an invariant from the guarded
fragment (Section~\ref{sec:calculus}) is buffer-free and holds of a memory
under any nonempty set of completeness marks, which is what allows the pool to be
read at full completeness and handed over with the acquirer's mark alone.

\section{The repaired calculus}
\label{sec:calculus}

The judgment is again $\spec{J}{P}{c}{Q}$. The rules \rn{skip},
\rn{assume}, \rn{assert}, \rn{fence}, \rn{store}, \rn{seq}, \rn{cons},
\rn{disj}, \rn{choice}, \rn{loop} and \rn{frame-sp} are those of
Figure~\ref{fig:original}. Figure~\ref{fig:repaired} shows the rest.

\begin{figure}[!ht]
\small
\begin{mathpar}
\inferrule[lock]{ }
  {\spec{J}{\femp}{\clock_i}{(\flock{i} \fsep J) \fseq \fbar{i}}}
\and
\inferrule[unlock]
  {A \fseq \fbar{i} \sentails J \fsep B}
  {\spec{J}{\flock{i} \fsep A}{\cunlock_i}{B}}
\and
\inferrule[frame-tm]
  {\spec{J}{P}{c}{Q} \\ \mod{c} \cap \fv{R} = \nil \\
   c \text{ contains no } \cunlock}
  {\spec{J}{R \fseq P}{c}{R \fseq Q}}
\and
\inferrule[load-any]
  {x \notin \fv{J,e,e',P}}
  {\spec{J}{e \fcell e' \fsseq P}{\cload{x}{e}_i}
         {(e \fcell e' \fsseq P) \conj x = e'}}
\and
\inferrule[store-any]{ }
  {\spec{J}{e \fcell e'' \fsseq P}{\cstore{e}{e'}_i}
         {(e \fcell e'' \fsseq P) \fseq e \fwrite{i} e'}}
\and
\inferrule[assign]
  {x \notin \fv{J}}
  {\spec{J}{P\subst{e}{x}}{\cassign{x}{e}_i}{P}}
\and
\inferrule[ex]
  {\spec{J}{P}{c}{Q} \\ x \notin \fv{J,c,Q}}
  {\spec{J}{\exists x \st P}{c}{Q}}
\end{mathpar}
\caption{Rules of the repaired calculus that differ from
Figure~\ref{fig:original}. \rn{load} is the original axiom with
$x \notin \fv{J,e,i,e',P}$, and \rn{conc} is the original rule with
$\mod{c_k} \cap \fv{J} = \nil$ added.}
\label{fig:repaired}
\end{figure}

\paragraph{The temporal frame rule.}
\rn{frame-tm} is the dissertation's rule with one more side condition: the
command must not contain $\cunlock$, anywhere. The reason is the semantic
one of Section~\ref{sec:model}. What the rules derive for a command that never releases is
valid in overlap mode, where overlapping history is tolerated; the
general \rn{unlock} rule is provided only in disjoint mode. The mechanized judgment carries the mode as an
index and has 21 rules, among them a temporal frame rule for overlap mode
and a rule weakening overlap to disjoint (Appendix~\ref{app:rules}). The
judgment displayed here is its disjoint mode, and the displayed
\rn{frame-tm} is derived in it: a derivation for a command without
$\cunlock$ can be replayed in overlap mode, framed, and weakened
(\rocq{owned_derive_temporal_frame_unlock_free}). The side condition is exact for the calculus: among the derivable
specifications, those with an overlap-mode derivation are the ones whose
command contains no $\cunlock$
(\rocq{owned_overlap_derivable_iff_unlock_free}). It is also
needed: \rn{frame-tm} applied to a valid specification of a release can
yield an invalid one (\rocq{owned_disjoint_temporal_frame_counterexample}).

\paragraph{Lock and unlock.}
\rn{lock} delivers the invariant together with the acquirer's barrier.
\rn{unlock} takes the whole of the thread's state $A$, notes that releasing
drains the thread's buffer, and asks for a split into something satisfying
$J$ and a remainder $B$. Because $A$ is arbitrary, nothing needs to be
framed around a release: whatever the thread wants to keep goes into $A$
and comes out in $B$.

\paragraph{Invariants.}
The soundness theorem requires $J$ to come from the fragment
\[
  H \bnfdef \femp \bnfbar e \fcell e' \bnfbar H \fsep H \bnfbar H \disj H
    \bnfbar \exists x \st H \bnfbar b \conj H .
\]
A pure formula may appear only alongside a heap formula. A bare $\bexpt$
is not buffer-free: it holds of a memory with pending writes, and an
invariant containing it could swallow them at release. Some restriction is
necessary: with the original $x \fpointsto 7$ as invariant, \rn{lock}
derives a specification that is false, because one acquisition cannot
supply every processor's mark
(\rocq{owned_soundness_without_fragment_refuted}).

\subsection{A proof system for assertions}
\label{sec:assertion-rules}

The dissertation left entailment between assertions semantic. The judgment
$P \vdash_a Q$ has 47 rules, none of which refers to satisfaction:
\begin{itemize}[nosep]
  \item the usual rules for $\conj$, $\disj$, $\exists$, $\forall$,
        substitution and equality, and a rule admitting a valid implication
        between boolean expressions;
  \item for $\fsep$ and $\fseq$: units, associativity, monotonicity,
        distribution over $\disj$ and $\exists$, movement of a pure
        conjunct in and out, and commutativity of $\fsep$;
  \item the three adjunctions for the implications;
  \item for barriers, locks and writes, the surviving laws of
        Appendix~\ref{app:statements};
  \item two laws for cells: cells at distinct addresses compose spatially
        exactly when they compose temporally, and
        $e \fcell u \fseq e \fcell v \vdash_a e \fcell v$.
\end{itemize}
Strong temporal conjunction has no rules of its own. It is an abbreviation,
and what is true of it follows from the rules for $\conj$, $\fsep$ and
$\fseq$. The system is sound (\rocq{assertion_derivation_sound}) for all
assertions and all processor domains. Exclusivity of cells, collapse of a
drained history onto its last write, and hiding of a cell's value are
derived in it, and the bound-variable renaming they need is itself proved. No completeness is claimed.
Appendix~\ref{app:assertion-rules} displays the rules that are particular
to this logic.

\paragraph{What became of the derived rules.}
The dissertation also lists a dozen derived and alternative rules. Most
are conveniences built from the primitive ones, and a few are proposals
for stronger axioms. \rn{load-p} and
\rn{store-p} are the load and store axioms restated for the all-processor
points-to $e \fpointsto e'$. The rules whose names end in \rn{-g}, such as
\rn{fence-g} and \rn{lock-g}, take an arbitrary precondition $P$ and append
the command's effect to it; for example, \rn{fence-g} is
$\spec{J}{P}{\cfence_i}{P \fseq \fbar{i}}$. Those ending in \rn{-b} are the
backward forms, which start from an arbitrary postcondition. \rn{locked} and \rn{atomic} give a specification to a body run between
acquiring and releasing the lock; \rn{atomic}, proposed as a stronger
axiom, lets the body rely on the barrier that the release supplies when it
restores the invariant.

In the repaired calculus \rn{fence-g} is derived exactly as the
dissertation displays it, and \rn{load-p} keeps its original form with the
freshness of the loaded variable made explicit. \rn{lock-g}, \rn{lock-b} and the
\rn{locked} wrapper are derived in a form that carries the transferred
invariant. The unlock variants are derived for the empty invariant, the
general case being \rn{unlock} itself. The ``daring'' rules for unlocked
access to shared state are unsound for the repaired calculus;
Section~\ref{sec:daring} gives a restricted form that is sound. Appendix~\ref{app:derived}
states each rule in its repaired form.

\section{Soundness and adequacy}
\label{sec:results}

\begin{theorem}[Soundness, \rocq{owned_calculus_sound}]
Let $J$ be in the invariant fragment and let there be two distinct active
processors. If $\spec{J}{P}{c}{Q}$ is derivable then
$\mtruespec{J}{\modedj}{P}{c}{Q}$, and if $c$ contains no $\cunlock$ then
also $\mtruespec{J}{\modeov}{P}{c}{Q}$.
\end{theorem}

The proof is by induction on the mode-indexed derivation, one semantic
lemma per rule; the overlap-mode half first replays the derivation in
overlap mode. Both hypotheses are needed: dropping either admits a
derivation of an invalid specification (Sections~\ref{sec:changes}
and~\ref{sec:calculus}).
No rule has a premise about safety or validity. \rn{cons} and \rn{unlock}
have entailments as premises, which may be discharged semantically or in
the proof system of Section~\ref{sec:assertion-rules}.

\begin{theorem}[Adequacy, \rocq{owned_derivation_adequacy}]
Under the same hypotheses, let $\spec{J}{P}{c}{Q}$ be derivable. Take a
machine memory with the lock free that splits spatially into a pool
satisfying $J$ and a resource satisfying $P$. Then no finite execution of
$c$ from that memory aborts. If an execution terminates, its final memory
splits into a pool and a resource satisfying $Q$, where either the lock is
free and the pool satisfies $J$, or the resource holds the lock and the
pool is empty.
\end{theorem}

Adequacy instantiates the external frame as empty, which satisfies both
modes, so a client never discharges a mode condition.

\paragraph{Share.}
The original \rn{share} is unsound next to the transferring \rn{lock}: a
command that ends with the lock held has the invariant in its own hands,
and \rn{share} would count it twice. It is sound when every
model of $P$ and of $Q$ is lock-free: then $\mtruespec{J}{m}{P}{c}{Q}$
implies $\mtruespec{\femp}{m}{J \fsep P}{c}{J \fsep Q}$, in either mode. The condition
must be semantic; a postcondition that merely omits the token syntactically
can still hide it under $\bexpt$.

For a specification from $\femp$ to $\femp$ there is a sharper statement: if the external frame is lock-free to begin with, every completed execution
ends with the lock free, and the final memory is exactly an invariant pool
together with a flushed descendant of that frame.

These are partial-correctness results. They say nothing of termination,
fairness or deadlock, and an execution stopped inside a critical section
need not satisfy the invariant.

The standing hypotheses are used less than one might expect. The frame
rules, the concurrency rule and the heap axioms need neither two processors
nor, except for \rn{frame-tm}, the fragment.

\section{How proofs proceed}
\label{sec:proofs}

The dissertation gave few examples. This section gives five, in increasing
size. All are derivations in the calculus of Section~\ref{sec:calculus},
checked by Rocq. Appendix~\ref{app:examples} gives the programs and
specifications of the larger ones in full.

\subsection{A store, with and without a fence}

For any invariant,
\[
  \spec{J}{x \fcell u}{\cseq{\cstore{x}{v}_i}{\cfence_i}}{x \fcell v}
\]
is derivable: \rn{store-any} gives
$x \fcell u \fseq x \fwrite{i} v$, \rn{frame-tm} carries that history
across $\cfence_i$, and a drained history collapses to its last write. This
is the store rule of ordinary separation logic. The same specification
without the fence is \emph{invalid}:
the store may terminate with its write still buffered, and a cell assertion
has no buffered writes. The logic is the difference between these two.

\subsection{Reading a buffered write}

After $\cstore{x}{v}_i$ the thread holds $R \fseq x \fwrite{i} v$, where
$R$ describes the older state of $x$. The load axiom wants the loaded write
to be spatially separate from the rest, and an old cell cannot be separated
spatially from a newer write to the same address. The temporal frame rule
is what resolves this. The newest write alone satisfies the axiom's
precondition, with $\femp$ for $P$; the older history is a
temporal frame, provided $r$ is free in none of $J$, $R$, $x$, $i$, $v$:
\[
\inferrule*[right=frame-tm]
  {\spec{J}{x \fwrite{i} v}{\cload{r}{x}_i}
         {x \fwrite{i} v \conj r = v}}
  {\spec{J}{R \fseq x \fwrite{i} v}{\cload{r}{x}_i}
         {R \fseq (x \fwrite{i} v \conj r = v)}}
\]
Store-to-load forwarding, the
behaviour most characteristic of the machine, is thus an instance of the
rule that most distinguishes the logic, and the reason it was worth
keeping.

\subsection{A shared container}

One cell $x$ is shared through a flag $f$. With
$X = \exists v \st x \fcell v$, the invariant is
\[
  J \;=\; f \fcell 0 \;\disj\; (f \fcell 1 \fsep X).
\]
A thread on processor $1$ takes the item, writes to it privately, and
returns it:
\begin{align*}
 & \clock_1;\ \cload{r}{f}_1;\ \cassume{r = 1}_1;\ \cstore{f}{0}_1;\
   \cunlock_1; \\
 & \cstore{x}{9}_1; \\
 & \clock_1;\ \cload{r}{f}_1;\ \cassume{r = 0}_1;\ \cstore{f}{1}_1;\
   \cunlock_1 .
\end{align*}
Meanwhile processor $2$ has an unrelated write $R = y \fwrite{2} 17$
buffered. The derived specification is
$\spec{J}{R}{c}{R}$. Writing
$H(a,v)$ for $(\exists u \st a \fcell u) \fseq a \fwrite{1} v$, the proof
passes through these assertions, each with $R$ as a spatial frame:
\begin{center}\small
\begin{tabular}{@{}ll@{}}
\toprule
After & The thread holds \\
\midrule
first load and assume & $\flock{1} \fsep (f \fcell 1 \fsep X)$ \\
$\cstore{f}{0}$ & $\flock{1} \fsep (H(f,0) \fsep X)$ \\
first unlock & $X$ \\
$\cstore{x}{9}$ & $H(x,9)$ \\
second lock & $H(x,9) \fseq ((\flock{1} \fsep J) \fseq \fbar{1})$ \\
second load and assume & $\flock{1} \fsep (f \fcell 0 \fsep x \fcell 9)$ \\
$\cstore{f}{1}$ & $\flock{1} \fsep (H(f,1) \fsep x \fcell 9)$ \\
second unlock & $\femp$ \\
\bottomrule
\end{tabular}
\end{center}
Three things are visible in it.

\emph{The foreign write costs one rule application.} $R$ concerns an
address the thread never touches, so a single \rn{frame-sp} around the
entire program carries it, across both releases.

\emph{The temporal frame rule carries the private write across the second
acquisition.} $H(x,9)$ is overlapping history, possibly still buffered.
\rn{frame-tm} places it before the acquired state, and the barrier that
\rn{lock} delivers then drains it, leaving $x \fcell 9$.

\emph{Releases are justified by entailment.} The two \rn{unlock} premises
are $(H(f,0) \fsep X) \fseq \fbar{1} \sentails J \fsep X$ and
$(H(f,1) \fsep x \fcell 9) \fseq \fbar{1} \sentails J \fsep \femp$. The
first keeps the item and the second donates it.

The example is not vacuous. From a memory with the lock free, $f = 1$,
$x = 7$ and processor $2$'s write pending, there is a complete 24-step
execution in which that write is pending in every state, and every execution from that memory is
safe and, if it completes, restores the invariant.

A second version replaces each $\cassume{}$ by a loop that releases the
lock and retries when the flag is wrong.
Its return loop has a branch for ``the container is already full'', and
that branch is derived from $\bexpf$: the thread holds $x$, so the pool
cannot.

\paragraph{The cost.}
The container's derivation applies \rn{cons} ten times and needs sixteen
entailments that are not identities. Most are rearrangements, a few are the
laws about barriers and histories used above, and the hardest is at the
second acquisition. All sixteen are derived in the assertion proof system
of Section~\ref{sec:assertion-rules}, and none of its 47 rules had to be
added to for the purpose. The same is true of the two examples that follow.
The retrying version of the container has not been migrated and still
discharges some entailments semantically.

\subsection{The dissertation's example}
\label{sec:own-example-repaired}

The message-passing program of the dissertation's Section~4.1 is
\[
  \cseq{\clock_1}{\cseq{\cstore{d}{1}_1}{\cunlock_1}}
  \;\oppar\;
  \cseq{\clock_2}{\cseq{\cload{x}{d}_2}{\cseq{\cunlock_2}
     {\cifthen{x = 1}{(\cseq{\cstore{d}{2}_2}{\cfence_2})}}}}
\]
and its claim is that from $d \fpointsto 0$ every terminating execution
ends with $d \fpointsto x + 1$. The original proof tracked progress in two
registers mentioned by the invariant. The repaired calculus forbids that,
and uses ownership instead. Give the sender a token cell $t$ and the
receiver a token cell $r$, neither ever read or written, and take
\[
  J \;=\; d \fcell 0 \;\disj\; (d \fcell 1 \fsep t \fcell 0) \;\disj\;
          (t \fcell 0 \fsep r \fcell 0).
\]
A sender still holding $t$ can only find the first disjunct when it
acquires the lock, so its store is safe; it donates $t$ with the updated
cell. A receiver holding $r$ cannot find the third. If it reads $1$ it
keeps $d$ and donates $r$, and then writes to $d$ with no lock. The
derivation uses no auxiliary code and no semantic entailment.

The tokens are then removed from the conclusion. An execution from the
original one-cell heap lifts to one with two extra untouched cells, and the
machine never changes the set of allocated addresses, which rules out the
one remaining disjunct in the $x = 0$ case. The result is a theorem about
the original program and the original heap
(\rocq{message_passing_all_runs}): no execution aborts, and every
terminating one ends with $x \in \set{0,1}$ and exactly the heap
$d \mapsto x + 1$, which gives the dissertation's
$d \fpointsto x + 1$.

\subsection{Named locks and a mailbox}

The machine has one lock. Programs want many. A named lock is a flag cell
$l$ with an invariant $I_l$, represented in the global invariant by
$(l \fcell 0 \fsep I_l) \disj l \fcell 1$. Acquiring it is a loop: take the
global lock, read the flag, set it if it was clear, release the global
lock, and retry otherwise. For any finite family of such locks, with $J$
the spatial product of their representations, the following are derived
rules, given a scratch register that does not occur in $J$:
\[
  \spec{J}{\femp}{\mathit{acquire}(l)}{I_l}
  \qquad\qquad
  \inferrule{I_l \fsep I_l \vdash_a \bexpf}
            {\spec{J}{I_l}{\mathit{release}(l)}{\femp}}
\]
The premise on release says that $I_l$ cannot be held twice, which is how a
thread holding $I_l$ knows the flag is set. This is the lock rule of
concurrent separation logic, recovered inside a logic with a single global
lock. Between a named acquire and its release the global lock is free.

The client is a mailbox guarded by a named lock. A producer writes a
message of $K$ cells with ordinary stores and no fence, takes the named
lock, marks the mailbox full, and releases. No fence is needed because the
machine will not let it take the global lock, inside the named acquire,
until its buffered writes have drained. A consumer retries until the mailbox
is full, takes the cells out, and reads them privately. The derived
specification says that the consumer ends up owning cells that hold the
values the producer wrote. It holds for every $K$, including zero, and by adequacy for every execution from
an actual initial memory, with the flags, the message cells and one further
address all distinct, in which a third processor's write to that address is
pending (\rocq{mailbox_all_runs}). A checked 95-step execution of the two-cell
instance leaves that write pending to the end.

\section{Unlocked access}
\label{sec:daring}

The dissertation opens with a simplification of Dekker's algorithm, which
a fence after each store makes correct on this machine:
\[
  \cseq{\cstore{f_1}{1}_1}{\cseq{\cfence_1}{\cload{r_1}{f_2}_1}}
  \;\oppar\;
  \cseq{\cstore{f_2}{1}_2}{\cseq{\cfence_2}{\cload{r_2}{f_1}_2}}
\]
From $f_1 = f_2 = 0$ the two registers should never both end at $0$. Each
thread reads a flag that the other writes, and neither takes the lock. The
calculus of Section~\ref{sec:calculus} cannot express this, because there
a thread may touch a cell only if it owns it or holds the lock.

\paragraph{The original rules.}
The dissertation's basic ``daring'' rules were meant for this case. They
let a thread load from or store to the invariant's cells without the lock,
provided that the single command, specified with $J$ in hand, gives $J$
back:
\[
  \inferrule{\spec{\femp}{J \fsep P}{a}{J \fsep Q}}{\spec{J}{P}{a}{Q}}
\]
The dissertation expected their soundness to be hard to show and left them
as future work. For the repaired calculus they are unsound. Adding the load rule lets
one derive a false specification
(\rocq{daring_load_extension_breaks_adequacy}), and each rule has an
instance whose premise is valid and whose conclusion is invalid
(\rocq{daring_load_honest_counterexample},
\rocq{daring_store_honest_counterexample}). Two features of the repaired
model are responsible. A thread that holds the lock also holds the
invariant's cells, so $J \fsep P$ can be false of a state the thread is
really in. And between lock operations the invariant's cells are fixed and
free of buffered writes, so a store without the lock can neither take the
cell it writes to nor leave its write there.

\paragraph{A restricted form.}
A second judgment, $\bspec{J}{P}{c}{Q}$, is defined for commands that
contain neither $\clock$ nor $\cunlock$. It differs from the first in
three ways.
\begin{itemize}[nosep]
  \item The invariant's resource may hold buffered writes, and $J$ may
    describe them with a history $e \fcell u \fseq e \fwrite{i} v$. A
    daring store places its write in the invariant. The dissertation noted
    that \rn{daring-store} requires this, and preferred invariants that
    describe the heap alone.
  \item Any primitive step may move cells between the invariant and the
    thread, where the first judgment allows this only at $\clock$ and
    $\cunlock$.
  \item Frames are disjoint from the thread's resource, and there is no
    temporal frame rule.
\end{itemize}
The rule for access is
\[
  \inferrule*[right=atomic]
    {\langle J \fsep P \rangle\; a\; \langle J \fsep Q \rangle}
    {\bspec{J}{P}{a}{Q}}
\]
for a single store, load, fence, assignment, $\cassume{}$ or
$\cassert{}$. Its premise is a statement about one machine step. From every
unlocked memory that splits into a model of $J$ and a model of $P$, the
step does not abort, and every result splits into a model of $J$ and a
model of $Q$. The invariant's part is read at full completeness, and the
thread's part gains no completeness mark except its own at a fence. No
condition is needed for the flushes that happen between steps, because
assertions are closed under flushing.

With \rn{atomic} the judgment has sixteen rules: the structural rules, the
spatial frame rule, the concurrency rule, and axioms for cells the thread
owns outright. It is sound and adequate for the machine
(\rocq{buffered_calculus_sound}, \rocq{buffered_derivation_adequacy}), with
no assumption on the number of processors.

\paragraph{The result.}
The program above is derivable with a postcondition that excludes
$r_1 = r_2 = 0$ (\rocq{buffered_dekker_derived};
Appendix~\ref{app:dekker} gives the invariant). Erasing the proof's token
cells, as for message passing, gives a theorem about the machine alone
(\rocq{dekker_fenced_all_runs}): from the heap $f_1 = f_2 = 0$ with empty
buffers, no execution aborts and no terminating execution ends with
$r_1 = 0$ and $r_2 = 0$. Without the fences such an execution exists
(\rocq{dekker_unfenced_zero_zero}).

\paragraph{Limits.}
This is the dissertation's simplified program, without the retry loop of
the full algorithm. The two
judgments are separate: a program cannot yet take the lock in one place and
use \rn{atomic} in another, and it is not known how to combine them. And a
thread needs a completeness mark before it can take a cell out of the
invariant, since a cell cannot be owned without one. A thread that starts
from $\femp$ may have none
(\rocq{buffered_emp_load_withdrawal_premise_invalid}). Starting each
thread with $\fbar{i}$, which costs nothing when the program begins,
supplies one (\rocq{buffered_barrier_parallel_initial_adequacy},
\rocq{buffered_marked_cell_withdrawal}).

\section{The mechanization}
\label{sec:mechanization}

The development is about 34,000 lines of Rocq in 218 files, using only the
standard library. It declares no axioms and admits no proofs. An audit file
prints the assumptions of 1,483 named results, and the build checks the
compiled proof objects independently. The proofs need neither
extensionality nor choice nor a finite set of processors. The
development also keeps, as a record, the repairs that were tried and set
aside.

The transcription of the original definitions is kept separate from every
repair, so that every claim in Section~\ref{sec:wrong} is a
theorem about the dissertation's own definitions. The report names only the principal definitions and theorems, at commit
\texttt{521acd6} of the development. The file
\texttt{paper/theorem-names.md} there lists the theorems behind its other
claims.

The work was done with AI assistance, in a loop that turned out to matter.
Two models, Claude Fable and ChatGPT Astra, took part: one developed the
proofs and the other reviewed them independently and proposed the design.
Each rule or lemma was stated before it was proved, and where a bounded
search was cheap it was run first. Twice the search or the proof attempt
refuted a confident argument on paper: a frame lemma that let the frame
evolve freely, and the soundness of \rn{frame-stm}. Both now appear above
as counterexamples. I do not think either would have been caught by
reading.

\section{Later work}
\label{sec:related}

This is a brief orientation, and it relies on the papers' own accounts of
their results.

\paragraph{Before and alongside.}
Ridge~\cite{ridge10} gave a rely-guarantee proof system for x86-TSO,
mechanized in HOL4. Cohen and Schirmer~\cite{cohen10} and
Owens~\cite{owens10} gave conditions on a program under which its TSO
executions are sequentially consistent, so that buffers need not be
reasoned about at all. The dissertation, and the proposal with Berdine
that preceded it~\cite{wehrman11}, sat between these: buffers are visible
in assertions, and separation is meant to keep the reasoning local.

\paragraph{iCAP-TSO.}
The closest successor is the logic of Sieczkowski, Svendsen, Birkedal and
Pichon-Pharabod~\cite{sieczkowski15}, which cites the dissertation. It
extends a higher-order concurrent separation logic with store buffers. Its
connectives include one stating that an assertion holds until a buffered
update is flushed and another afterwards, which are temporal connectives
in the dissertation's sense, attached to a single update. Its central idea
is a \emph{fiction of sequential consistency}: a library that synchronizes
enough gets a specification from which its clients reason in ordinary
separation logic. The dissertation's logic has nothing comparable. They
also verify a spin lock whose release is a plain buffered store, which is
still out of reach here (Section~\ref{sec:open}). Jacobs~\cite{jacobs14}
proposed ``TSO spaces'' for separation logic with a related aim.

\paragraph{Elsewhere.}
Bornat, Alglave and Parkinson~\cite{bornat15} built a rely-guarantee logic
for weak memory around temporal modalities such as ``since'', so the idea
that ordering between writes belongs in the assertion language has
reappeared in another form. Most later work targets the C11 model and
weaker ones: relaxed separation logic~\cite{vafeiadis13} and
GPS~\cite{turon14}, and their reconstruction inside Iris~\cite{kaiser17},
which made mechanized soundness proofs the norm for such logics. For x86,
Pierogi~\cite{bila22} is an Owicki--Gries logic over a view-based
semantics that also covers persistent memory.

\paragraph{Where this report stands.}
Against that background, what the repaired logic offers is narrow and, as
far as I know, still distinctive: sequential and spatial separating
conjunctions over the whole history of a thread's writes, each with a frame
rule, and a soundness proof checked end to end against an operational TSO
machine. Section~\ref{sec:trilemma}'s conflict between temporal framing,
units and ownership transfer does not depend on TSO and may be of use to
anyone designing a logic with a sequential conjunction.
Section~\ref{sec:daring} handles unlocked access to shared state in a
separate judgment. It remains open whether a single judgment for this
logic can handle both locked and unlocked access in the same program
(Section~\ref{sec:open}). A further question is how to hide the memory
model from library clients, as iCAP-TSO does.

\section{What remains open}
\label{sec:open}

\paragraph{Unlocked access to shared state.}
The dissertation opens with Dekker's algorithm, made correct on this
machine by a fence after each store. Each thread reads a flag that the
other writes, with no lock. The calculus cannot express that: a thread may
touch a cell only if it owns it or holds the lock. The dissertation's basic ``daring'' rules were meant for this. They would let a
thread load from or store to the invariant's cells without the lock,
provided that the single command, specified with $J$ in hand, gives $J$
back. The dissertation doubted their soundness and left them as future
work. For the repaired calculus the doubt was justified. Adding the load
rule lets one derive a false specification
(\rocq{daring_load_extension_breaks_adequacy}), and each rule has an
instance whose premise is valid and whose conclusion is invalid
(\rocq{daring_load_honest_counterexample},
\rocq{daring_store_honest_counterexample}). The store rule fails even
while the lock is free. A sound rule for unlocked access, and with it a
proof of Dekker's algorithm, is still missing.

\paragraph{Realistic locks.}
The one lock is global and blocks every other processor. Programmer-level
locks can be encoded as flag cells under the invariant, as the container
does, but each release must itself take the global lock. A spinlock
released by a bare store is another instance of unlocked access.

\paragraph{A strong temporal connective.}
The frame property holds for a single split that is both spatial and
temporal (\rocq{owned_safe_common_frame}). Making that a connective of the
assertion language, in place of the conjunction of two, would restore
\rn{frame-stm}. Whether its other laws would then hold is not known.

\paragraph{Assertions.}
The proof system for assertions is sound and was sufficient for every
derivation attempted. Nothing is known about what it cannot prove, and
nothing is automated.

\paragraph{The rest of the dissertation.}
Allocation, recursive predicates, permissions and the other topics of
Chapter~5 are untouched, as is the sequential logic of Chapter~3.


\begin{thebibliography}{99}
\bibitem{bila22}
E.~V. Bila, B.~Dongol, O.~Lahav, A.~Raad and J.~Wickerson.
\newblock View-based Owicki--Gries reasoning for persistent x86-TSO.
\newblock In \emph{ESOP}, 2022.
\bibitem{bornat15}
R.~Bornat, J.~Alglave and M.~Parkinson.
\newblock New lace and arsenic: adventures in weak memory with a program
  logic.
\newblock arXiv:1512.01416, 2015.
\bibitem{cohen10}
E.~Cohen and B.~Schirmer.
\newblock From total store order to sequential consistency: a practical
  reduction theorem.
\newblock In \emph{ITP}, 2010.
\bibitem{jacobs14}
B.~Jacobs.
\newblock Verifying TSO programs.
\newblock Report CW660, KU Leuven, 2014.
\bibitem{kaiser17}
J.-O. Kaiser, H.-H. Dang, D.~Dreyer, O.~Lahav and V.~Vafeiadis.
\newblock Strong logic for weak memory: reasoning about release-acquire
  consistency in Iris.
\newblock In \emph{ECOOP}, 2017.
\bibitem{owens10}
S.~Owens.
\newblock Reasoning about the implementation of concurrency abstractions on
  x86-TSO.
\newblock In \emph{ECOOP}, 2010.
\bibitem{ridge10}
T.~Ridge.
\newblock A rely-guarantee proof system for x86-TSO.
\newblock In \emph{VSTTE}, 2010.
\bibitem{sieczkowski15}
F.~Sieczkowski, K.~Svendsen, L.~Birkedal and J.~Pichon-Pharabod.
\newblock A separation logic for fictional sequential consistency.
\newblock In \emph{ESOP}, 2015.
\bibitem{turon14}
A.~Turon, V.~Vafeiadis and D.~Dreyer.
\newblock GPS: navigating weak memory with ghosts, protocols, and
  separation.
\newblock In \emph{OOPSLA}, 2014.
\bibitem{vafeiadis13}
V.~Vafeiadis and C.~Narayan.
\newblock Relaxed separation logic: a program logic for C11 concurrency.
\newblock In \emph{OOPSLA}, 2013.
\bibitem{wehrman11}
I.~Wehrman and J.~Berdine.
\newblock A proposal for weak-memory local reasoning.
\newblock In \emph{LOLA}, 2011.
\bibitem{wehrman12}
I.~Wehrman.
\newblock \emph{Weak-Memory Local Reasoning}.
\newblock PhD dissertation, The University of Texas at Austin, December 2012.
\end{thebibliography}

\appendix
\section{The repaired calculus in full}
\label{app:rules}

Figure~\ref{fig:full} gives the mechanized calculus, whose judgment
$\mspec{J}{m}{P}{c}{Q}$ carries a mode $m$, either $\modeov$ or $\modedj$.
The judgment of Section~\ref{sec:calculus} is its disjoint mode.

\begin{figure}[!ht]
\small
\begin{mathpar}
\inferrule[skip]{ }{\mspec{J}{m}{P}{\cskip_i}{P}}
\and
\inferrule[assume]{ }{\mspec{J}{m}{\neg b \disj P}{\cassume{b}_i}{P}}
\and
\inferrule[assert]{ }{\mspec{J}{m}{b \conj P}{\cassert{b}_i}{P}}
\and
\inferrule[assign]{x \notin \fv{J}}
  {\mspec{J}{m}{P\subst{e}{x}}{\cassign{x}{e}_i}{P}}
\and
\inferrule[fence]{ }{\mspec{J}{m}{\femp}{\cfence_i}{\fbar{i}}}
\and
\inferrule[lock]{ }
  {\mspec{J}{m}{\femp}{\clock_i}{(\flock{i} \fsep J) \fseq \fbar{i}}}
\and
\inferrule[unlock]{A \fseq \fbar{i} \sentails J \fsep B}
  {\mspec{J}{\modedj}{\flock{i} \fsep A}{\cunlock_i}{B}}
\and
\inferrule[load]{x \notin \fv{J,e,i,e',P}}
  {\mspec{J}{m}{e \fwrite{i} e' \fsseq P}{\cload{x}{e}_i}
         {(e \fwrite{i} e' \fsseq P) \conj x = e'}}
\and
\inferrule[store]{ }
  {\mspec{J}{m}{e \fwrite{i} e'' \fsseq P}{\cstore{e}{e'}_i}
         {(e \fwrite{i} e'' \fsseq P) \fseq e \fwrite{i} e'}}
\and
\inferrule[load-any]{x \notin \fv{J,e,e',P}}
  {\mspec{J}{m}{e \fcell e' \fsseq P}{\cload{x}{e}_i}
         {(e \fcell e' \fsseq P) \conj x = e'}}
\and
\inferrule[store-any]{ }
  {\mspec{J}{m}{e \fcell e'' \fsseq P}{\cstore{e}{e'}_i}
         {(e \fcell e'' \fsseq P) \fseq e \fwrite{i} e'}}
\and
\inferrule[frame-sp]
  {\mspec{J}{m}{P}{c}{Q} \\ \mod{c} \cap \fv{R} = \nil}
  {\mspec{J}{m}{R \fsep P}{c}{R \fsep Q}}
\and
\inferrule[frame-tm]
  {\mspec{J}{\modeov}{P}{c}{Q} \\ \mod{c} \cap \fv{R} = \nil}
  {\mspec{J}{\modeov}{R \fseq P}{c}{R \fseq Q}}
\and
\inferrule[weaken]{\mspec{J}{\modeov}{P}{c}{Q}}{\mspec{J}{\modedj}{P}{c}{Q}}
\and
\inferrule[cons]
  {P \sentails P' \\ \mspec{J}{m}{P'}{c}{Q'} \\ Q' \sentails Q}
  {\mspec{J}{m}{P}{c}{Q}}
\and
\inferrule[seq]
  {\mspec{J}{m}{P}{c}{R} \\ \mspec{J}{m}{R}{c'}{Q}}
  {\mspec{J}{m}{P}{\cseq{c}{c'}}{Q}}
\and
\inferrule[choice]
  {\mspec{J}{m}{P}{c}{Q} \\ \mspec{J}{m}{P}{c'}{Q}}
  {\mspec{J}{m}{P}{\cchoice{c}{c'}}{Q}}
\and
\inferrule[loop]{\mspec{J}{m}{P}{c}{P}}{\mspec{J}{m}{P}{\cloop{c}}{P}}
\and
\inferrule[disj]
  {\mspec{J}{m}{P}{c}{Q} \\ \mspec{J}{m}{P'}{c}{Q}}
  {\mspec{J}{m}{P \disj P'}{c}{Q}}
\and
\inferrule[ex]
  {\mspec{J}{m}{P}{c}{Q} \\ x \notin \fv{J,c,Q}}
  {\mspec{J}{m}{\exists x \st P}{c}{Q}}
\and
\inferrule[conc]
  {\mspec{J}{m}{P_k}{c_k}{Q_k} \\
   \mod{c_{1-k}} \cap \fv{J,P_k,c_k,Q_k} = \nil \quad (k = 0,1)}
  {\mspec{J}{m}{P_0 \fsep P_1}{\cpar{c_0}{c_1}}{Q_0 \fsep Q_1}}
\end{mathpar}
\caption{The 21 rules of the repaired calculus (\rocq{owned_derivation}).
In \rn{lock} and \rn{unlock} the processor is a constant.}
\label{fig:full}
\end{figure}

\FloatBarrier
\section{The safety predicate in full}
\label{app:safety}

This appendix spells out the definitions that Section~\ref{sec:model} gives
in prose (\rocq{owned_safe}, \rocq{owned_valid}).

\paragraph{Resources and futures.}
Every memory here blocks only active processors and has empty buffers at
inactive ones. A \emph{resource} $\nu = (\mu,\Gamma)$ is a generalized
memory whose marks $\Gamma$ are active processors, and whose heap is empty
if $\Gamma$ is. A memory \emph{holds the lock for $i$} when $i$ is active
and the memory blocks exactly the other active processors. Its \emph{futures}, written $\nu' \modrefines \nu$, are the
resources reachable by any number of steps of two kinds:
\begin{itemize}[nosep]
  \item add marks: $(\mu,\Gamma)$ becomes $(\mu,\Gamma')$ with
        $\Gamma \subseteq \Gamma'$;
  \item flush: $(\mu,\Gamma)$ becomes $(\mu',\Gamma')$, where $\mu'$ is
        $\mu$ after the oldest buffered write of some active processor $i$
        not blocked in $\mu$ moves to the heap, and
        $\Gamma \cup \set{i} \subseteq \Gamma'$.
\end{itemize}

\paragraph{Environments.}
An \emph{environment} of $\nu$ at stack $s$, in mode $m$, consists of a
pool $\mu_J$, a frame resource $\nu_F$, a machine memory $\mu_0$ and a
phase $\phi$, which is either \textsf{free} or $\textsf{held}(i)$, such
that
\begin{enumerate}[nosep]
  \item $\mu_0 \in \mu_J \ssep \mu$ for the memory $\mu$ of some resource in
        $\nu_F \mhash \nu$;
  \item in mode $\modedj$, the addresses of $\nu_F$ and of $\nu$ are
        disjoint;
  \item if $\phi = \textsf{free}$, no processor is blocked in $\mu_0$ and
        $(s,\mu_J,\setprocessors) \sentails J$; and
  \item if $\phi = \textsf{held}(i)$, $\mu_0$ holds the lock for $i$, and
        $\mu_J$ is empty.
\end{enumerate}

\paragraph{Updates.}
A program step carries an action label $\alpha$: fence by $i$, acquire by
$i$, release by $i$, or ordinary. Stores, loads, assignments and the
structural steps of choice, iteration and composition are ordinary. Let
such a step take the local resource from $\nu = (\mu,\Gamma)$ to $\nu' = (\mu',\Gamma')$, the pool from
$\mu_J$ to $\mu'_J$ and the phase from $\phi$ to $\phi'$. The step is a
permitted \emph{update} when
\begin{enumerate}[nosep]
  \item \emph{transfer:}
    if $\alpha$ is an acquire by $i$ then $\phi = \textsf{free}$,
    $\phi' = \textsf{held}(i)$, $\mu'_J$ is empty and $\mu'$ holds the lock
    for $i$;
    if $\alpha$ is a release by $i$ then $\phi = \textsf{held}(i)$,
    $\phi' = \textsf{free}$, $\mu$ holds the lock for $i$ and $\mu'$ blocks
    no processor;
    otherwise $\phi' = \phi$, $\mu'_J = \mu_J$, and $\mu'$ blocks the same
    processors as $\mu$;
  \item \emph{footprint:} every address of $\mu'_J$ or $\mu'$ is an address
    of $\mu_J$ or $\mu$; and
  \item \emph{mark budget:} $\Gamma' \subseteq \Gamma \cup \set{i}$ if
    $\alpha$ is a fence, acquire or release by $i$, and
    $\Gamma' \subseteq \Gamma$ if $\alpha$ is ordinary.
\end{enumerate}

\paragraph{Safety and validity.}
$\safe{0}{c,s,\nu,J,Q}$ always holds. $\safe{n+1}{c,s,\nu,J,Q}$ holds in
mode $m$ when, for every future $\nu_1 \modrefines \nu$:
\begin{enumerate}[nosep]
  \item if $c$ is $\cskip$ then $s,\nu_1 \sentails Q$; and
  \item for every environment $(\mu_J,\nu_F,\mu_0,\phi)$ of $\nu_1$ at $s$,
    the configuration $c,(s,\mu_0)$ does not abort, and for every program
    step from it to $c',(s',\mu_1)$ there are $\nu'$, $\mu'_J$ and $\phi'$
    such that $(\mu'_J,\nu_F,\mu_1,\phi')$ is an environment of $\nu'$ at
    $s'$, the step is a permitted update from $(\nu_1,\mu_J,\phi)$ to
    $(\nu',\mu'_J,\phi')$, and $\safe{n}{c',s',\nu',J,Q}$.
\end{enumerate}
$\mtruespec{J}{m}{P}{c}{Q}$ holds when $s,\nu \sentails P$ implies
$\safe{n}{c,s,\nu,J,Q}$ in mode $m$ for every $n$. Flushes by the machine are not program steps. When $J$ is in the
invariant fragment the pool has no buffered writes, so each flush belongs
either to the local resource or to the frame
(\rocq{pooled_resource_flush_attribution}). A local flush is covered by
the quantification over futures. A flush of the frame is covered by
choosing the environment again, since the definition quantifies over all
of them.

\section{The repaired definitions}
\label{app:definitions}

Three definitions of the dissertation's model were replaced
(Table~\ref{tab:local}). The first two admit fewer compositions than the originals, and the third
adds a requirement to each implication.

\paragraph{Buffer overriding.}
Originally, $b \override b'$ is the set of interleavings of the buffers $b$
and $b'$ in which the most recent write to each address is the one that
$b'$ overriding $b$ would give. That constrains only the last write to each
address. In the repair, an interleaving is built from the oldest write
onwards, and a write to address $\ell$ may be taken from the later buffer
$b'$ only when no write to $\ell$ remains in $b$
(\rocq{ordered_overrides}). All of $b$'s writes to an address therefore
precede all of $b'$'s.

\paragraph{Temporal compatibility.}
Originally, $\mu_0$ may precede $\mu_1$ when three conditions hold. Their
blocked sets are disjoint. For each active processor $i$ \emph{unblocked
in both}, buffer $i$ of $\mu_0$ writes to no address in the heap of
$\mu_1$. For each pair of distinct active processors $i$ and $j$
\emph{unblocked in both}, buffer $i$ of $\mu_0$ and buffer $j$ of $\mu_1$
write to disjoint sets of addresses. The repair drops ``unblocked in
both'' from the second and third conditions (\rocq{stable_compatible}).

\paragraph{Implications.}
Originally, $\nu$ satisfies $P \fsepi Q$ when every composition of $\nu$
with a resource satisfying $P$ satisfies $Q$, and similarly for the two
temporal implications. In the repair the same is required of every future of $\nu$, so that an
implication which holds now cannot be falsified by a flush. This applies
at every implication in an assertion, nested ones included. It does not
follow that a repaired assertion has fewer models than the original one,
because the antecedent of an implication is reinterpreted too.

\FloatBarrier
\section{The dissertation's assertion laws}
\label{app:statements}

The dissertation states its laws for assertions in two figures, one of
equivalences and one of entailments (Figures~4.5 and~4.6).
Table~\ref{tab:algebra} classifies every one of them in the final
interpretation. A proved law holds for all
assertions, implications included, and all processor domains. Each refuted
law is also refuted in the dissertation's original interpretation.

\begin{table}[!ht]
\small
\begin{tabularx}{\textwidth}{@{}p{0.42\textwidth}lX@{}}
\toprule
Law ($\bullet$ ranges over $\fsep$, $\fseq$, $\fsseq$) & Status & Note \\
\midrule
$P \fgen \femp \sequiv P \sequiv \femp \fgen P$ & proved & all three \\
$(P \fgen Q) \fgen R \sequiv P \fgen (Q \fgen R)$
  & proved for $\fsep$, $\fseq$ & refuted for $\fsseq$, both directions \\
$P \fsep Q \sequiv Q \fsep P$ & proved & \\
$P \fgen \flock{e} \sequiv \flock{e} \fgen P$ & proved
  & indeed $P \fseq \flock{e} \sequiv P \fsep \flock{e}$ \\
$\flock{e} \fgen \flock{f} \sequiv \bexpf$ & qualified
  & needs a processor blocked by both \\
$\fbar{e} \fgen \fbar{e} \sequiv \fbar{e}$ & proved & \\
$(P \fgen Q) \fseq \fbar{e} \sequiv
 (P \fseq \fbar{e}) \fgen (Q \fseq \fbar{e})$ & proved & all three \\
$P \fsseq \fbar{e} \sequiv P \fseq \fbar{e}$ & proved & \\
$\flock{e} \fseq \fbar{e} \sequiv \flock{e} \fsep \fbar{e}$ & proved & \\
$e \fwrite{e'} e'' \fsep e \fwrite{f'} f'' \sequiv \bexpf$ & proved & \\
$e \fwrite{i} e' \fseq e \fwrite{i} f \fseq \fbar{i} \sequiv
 e \fwrite{i} f \fseq \fbar{i}$ & proved & \\
\midrule
$\fbar{e} \sentails \femp$; monotonicity of $\fgen$ & proved & \\
$e \fpointsto f \sentails e \fwrite{e'} f$ & qualified
  & $e'$ must denote an active processor \\
$P \fsseq Q \sentails P \fsep Q$ and $P \fsseq Q \sentails P \fseq Q$
  & proved & \\
Small exchange, $\fsseq$ outermost & refuted
  & under either reading of its premise \\
Small exchange between $\fsep$ and $\fseq$ & refuted or vacuous
  & refuted for the premise as displayed; read as an inclusion of
    connectives, the premise forces an empty processor set \\
Full exchange for $\fsep$ and $\fsseq$ & refuted & \\
\bottomrule
\end{tabularx}
\caption{The original assertion laws, classified.}
\label{tab:algebra}
\end{table}

The pattern is simple to state. The laws that rebracket or exchange
through $\fsseq$ fail, for the reason \rn{frame-stm} fails: the two
conjuncts of $P \fsseq Q$ choose their splits independently. No exchange
law between different conjunctions is usable. The remaining laws survive,
two of them with a side condition. The unit laws needed the
most care, since composing with an empty resource can add completeness
marks, and the proof must show that no assertion, however many implications
it contains, can tell.

\FloatBarrier
\section{A selection of the assertion rules}
\label{app:assertion-rules}

Figure~\ref{fig:assertion-rules} shows the rules of $P \vdash_a Q$ that are
particular to this logic. The figure omits the usual structural and logical rules, substitution and
equality, a rule admitting valid implications between boolean expressions
and a case split on a boolean, the distribution of $\fsep$ and $\fseq$
over $\disj$ and $\exists$, and the mirror image of the rule for pure
conjuncts. In the figure $\fgen$
stands for either $\fsep$ or $\fseq$, the same one throughout a rule, and
$P \dashv\vdash_a Q$ abbreviates the two rules $P \vdash_a Q$ and
$Q \vdash_a P$.

\begin{figure}[!ht]
\small
\begin{mathpar}
\femp \fgen P \dashv\vdash_a P \dashv\vdash_a P \fgen \femp
\and
(P \fgen Q) \fgen R \dashv\vdash_a P \fgen (Q \fgen R)
\and
P \fsep Q \vdash_a Q \fsep P
\and
\inferrule{P \vdash_a P' \\ Q \vdash_a Q'}{P \fgen Q \vdash_a P' \fgen Q'}
\and
(b \conj P) \fgen Q \dashv\vdash_a b \conj (P \fgen Q)
\and
\inferrule{P \fsep Q \vdash_a R}{Q \vdash_a P \fsepi R}
\and
\inferrule{P \fseq Q \vdash_a R}{Q \vdash_a P \fseqli R}
\and
\inferrule{Q \fseq P \vdash_a R}{Q \vdash_a R \fseqri P}
\\
\fbar{e} \vdash_a \femp
\and
\fbar{e} \fgen \fbar{e} \dashv\vdash_a \fbar{e}
\and
P \fseq \fbar{e} \vdash_a P \fsep \fbar{e}
\and
(P \fgen Q) \fseq \fbar{e} \dashv\vdash_a
  (P \fseq \fbar{e}) \fgen (Q \fseq \fbar{e})
\\
P \fseq \flock{e} \dashv\vdash_a P \fsep \flock{e}
\and
\flock{e} \fseq P \dashv\vdash_a \flock{e} \fsep P
\and
\inferrule{k \text{ active} \\ k \neq i \\ k \neq j}
  {\flock{i} \fsep \flock{j} \vdash_a \bexpf}
\\
e \fwrite{f} e' \fsep e \fwrite{f'} e'' \vdash_a \bexpf
\and
(e \fwrite{f} e' \fseq e \fwrite{f} e'') \fseq \fbar{f}
  \dashv\vdash_a e \fwrite{f} e'' \fseq \fbar{f}
\and
\inferrule{f \text{ denotes an active processor in every stack}}
  {e \fpointsto e' \vdash_a e \fwrite{f} e'}
\\
e \neq f \conj (e \fcell e' \fseq f \fcell f') \dashv\vdash_a
e \neq f \conj (e \fcell e' \fsep f \fcell f')
\and
e \fcell e' \fseq e \fcell e'' \vdash_a e \fcell e''
\end{mathpar}
\caption{Assertion rules particular to this logic, from the 47 of
\rocq{assertion_derivation}. Each implication rule is invertible; only the
introduction direction is shown.}
\label{fig:assertion-rules}
\end{figure}

A few of these do most of the work in the derivations of
Section~\ref{sec:proofs}. The barrier rules say that a barrier is an empty
resource that can be duplicated, and that draining a composite drains each
part. The lock rules say that a lock assertion carries no history, so its
temporal and spatial compositions agree. The collapse rule for writes says
that once a processor's writes in a history have drained, only its last
write to an address is visible. The two cell rules are what make cells behave like
ordinary heap cells: cells at different addresses can be reordered freely,
and a later cell at the same address overwrites an earlier one.

\FloatBarrier
\section{The dissertation's derived rules}
\label{app:derived}

Most of the dissertation's derived rules are conveniences built from the
primitive ones. Figure~\ref{fig:derived} states them as they are derived in the
repaired calculus, and Table~\ref{tab:derived} records how each differs
from its original display. \rn{load-p} and \rn{store-p} are the load and
store axioms for the all-processor points-to $e \fpointsto e'$, so that a
thread can use a cell without naming the processor that wrote it. The
\rn{-g} rules run forwards from an arbitrary precondition $P$ and append
the command's effect to it. The \rn{-b} rules run backwards from an
arbitrary postcondition $Q$, and their precondition says that $Q$ will hold
once the command's effect is appended. \rn{locked} gives a specification to
a body run between acquiring and releasing the lock.

\begin{figure}[!ht]
\small
\begin{mathpar}
\inferrule[load-p]{x \notin \fv{J,e,e',P}}
  {\spec{J}{e \fpointsto e' \fsseq P}{\cload{x}{e}_i}
        {(e \fpointsto e' \fsseq P) \conj x = e'}}
\and
\inferrule[store-p]{ }
  {\spec{J}{e \fpointsto e'' \fsseq P}{\cstore{e}{e'}_i}
        {(e \fpointsto e'' \fsseq P) \fseq e \fwrite{i} e'}}
\and
\inferrule[fence-g]{ }{\spec{J}{P}{\cfence_i}{P \fseq \fbar{i}}}
\and
\inferrule[fence-b]{ }{\spec{J}{Q \fseqri \fbar{i}}{\cfence_i}{Q}}
\and
\inferrule[lock-g]{ }
  {\spec{J}{P}{\clock_i}{P \fseq ((\flock{i} \fsep J) \fseq \fbar{i})}}
\and
\inferrule[lock-b]{ }
  {\spec{J}{Q \fseqri ((\flock{i} \fsep J) \fseq \fbar{i})}{\clock_i}{Q}}
\and
\inferrule[unlock-g]{ }
  {\spec{\femp}{P \fsep \flock{i}}{\cunlock_i}{P \fseq \fbar{i}}}
\and
\inferrule[unlock-b]{ }
  {\spec{\femp}{(Q \fseqri \fbar{i}) \fsep \flock{i}}{\cunlock_i}{Q}}
\and
\inferrule[locked]
  {\spec{J}{P \fseq ((\flock{i} \fsep J) \fseq \fbar{i})}{c}
        {\flock{i} \fsep A} \\
   A \fseq \fbar{i} \sentails J \fsep Q}
  {\spec{J}{P}{\cseq{\clock_i}{\cseq{c}{\cunlock_i}}}{Q}}
\end{mathpar}
\caption{The dissertation's derived rules in the form derived in the
repaired calculus. In the lock rules the processor is a constant.}
\label{fig:derived}
\end{figure}

\begin{table}[!ht]
\small
\begin{tabularx}{\textwidth}{@{}lX@{}}
\toprule
Original rule & In the repaired calculus \\
\midrule
\rn{load-p} & Derived with the original triple and the original
  points-to; the freshness of $x$ is made explicit. \\
\rn{store-p} & The display has a load command and a strong temporal
  suffix, which no memory satisfies; with a store command and a temporal
  suffix, as in \rn{store}, it is derived. \\
Strengthened \rn{lock}, \rn{unlock} & Subsumed: \rn{lock} already delivers
  $\fbar{i}$, and \rn{unlock}'s premise already assumes it. \\
\rn{fence-g} & Derived as displayed. \\
\rn{lock-g}, \rn{lock-b} & Derived with the transferred invariant made
  explicit. \\
\rn{unlock-g}, \rn{unlock-b} & Derived for the empty invariant; in general
  the thread cannot promise to keep all of $P$ while donating part of it. \\
\rn{fence-b} & The display has $\clock$ where $\cfence$ is meant; with
  $\cfence$ it is derived. \\
\rn{locked}, \rn{atomic} & Derived as a wrapper whose body is proved under
  the same $J$. Importing a body proved under $\femp$ is not recovered. \\
Daring load and store & Unsound for the repaired calculus. A restricted
  form is sound in a separate judgment (Section~\ref{sec:daring}). \\
\bottomrule
\end{tabularx}
\caption{Status of the rules in the dissertation's ``Derived and
Alternative Axioms and Inference Rules''.}
\label{tab:derived}
\end{table}

\FloatBarrier
\section{The programs of Section~\ref{sec:proofs}}
\label{app:examples}

This appendix gives the programs and specifications that
Section~\ref{sec:proofs} describes in prose. Conditionals and loops are the
dissertation's abbreviations for choice, iteration and $\cassume{}$.

\paragraph{The retrying container.}
With $J = f \fcell 0 \disj (f \fcell 1 \fsep X)$ as in
Section~\ref{sec:proofs}, taking the item is
\begin{align*}
\mathit{take} \;=\;\;
 & \cassign{g}{0}_1;\ \opwhile\ g = 0\ \opdo \\
 & \quad (\clock_1;\ \cload{r}{f}_1;\
     \opif\ r = 1\ \opthen\ (\cstore{f}{0}_1;\ \cunlock_1;\ \cassign{g}{1}_1)\
     \opelse\ (\cunlock_1;\ \cassign{g}{0}_1))
\end{align*}
and $\mathit{put}$ is the same with the test $r = 0$ and the store
$\cstore{f}{1}_1$. The program is
$\mathit{take};\ \cstore{x}{9}_1;\ \mathit{put}$, and the derived
specification is $\spec{J}{R}{c}{R}$ with $R = y \fwrite{2} 17$
(\rocq{owned_retrying_container_derived}).

\paragraph{Message passing.}
With $J$ as in Section~\ref{sec:own-example-repaired}, the two threads and
their composition have the specifications
\begin{align*}
 & \spec{J}{t \fcell 0}
    {\cseq{\clock_1}{\cseq{\cstore{d}{1}_1}{\cunlock_1}}}{\femp} \\
 & \spec{J}{r \fcell 0}{\mathit{receiver}}
    {(r \fcell 0 \conj x = 0) \disj (d \fcell 2 \conj x = 1)} \\
 & \spec{J}{t \fcell 0 \fsep r \fcell 0}{\mathit{sender} \oppar
    \mathit{receiver}}
    {(r \fcell 0 \conj x = 0) \disj (d \fcell 2 \conj x = 1)}
\end{align*}
(\rocq{mp_sender_derived}, \rocq{mp_receiver_derived},
\rocq{mp_parallel_derived}). The sender gives up its token. The receiver
ends either still holding its token, having read $0$, or owning the data
cell with the value $2$, having read $1$.

\paragraph{Named locks.}
For a flag cell $l$, a processor $i$ and a scratch register $x$:
\begin{align*}
\mathit{acquire}_i(l,x) \;=\;\;
 & \cassign{x}{1}_i;\ \opwhile\ x = 1\ \opdo \\
 & \quad (\clock_i;\ \cload{x}{l}_i;\
     \opif\ x = 0\ \opthen\ (\cstore{l}{1}_i;\ \cunlock_i)\
     \opelse\ \cunlock_i) \\
\mathit{release}_i(l) \;=\;\;
 & \clock_i;\ \cstore{l}{0}_i;\ \cunlock_i
\end{align*}
Under the invariant below the flag is always $0$ or $1$, so the loop exits
with $x = 0$, the value read from a flag that was clear and is now set. For locks $l_1,\dots,l_n$ with invariants $I_1,\dots,I_n$, the
global invariant is
\[
  J \;=\; \bigl((l_1 \fcell 0 \fsep I_1) \disj l_1 \fcell 1\bigr) \fsep
          \cdots \fsep
          \bigl((l_n \fcell 0 \fsep I_n) \disj l_n \fcell 1\bigr),
\]
and the two derived rules of Section~\ref{sec:proofs} hold for each $l_k$,
given $x \notin \fv{J}$.

\paragraph{The mailbox.}
There is one named lock $l$, a flag cell $m$ that says whether the mailbox
is full, and message cells $a_1,\dots,a_K$ that are to receive the values
$v_1,\dots,v_K$. Write $A(\vec{w})$ for
$a_1 \fcell w_1 \fsep \cdots \fsep a_K \fcell w_K$, which is $\femp$ when
$K = 0$. The lock's invariant is
\[
  I_l \;=\; m \fcell 0 \;\disj\; (m \fcell 1 \fsep A(\vec{v})),
\]
and $J = (l \fcell 0 \fsep I_l) \disj l \fcell 1$. The producer runs on
processor $1$ and the consumer on processor $2$; the processor
annotations on their stores, loads and assignments are omitted here:
\begin{align*}
\mathit{producer} \;=\;\;
 & \cstore{a_1}{v_1};\ \dots;\ \cstore{a_K}{v_K};\
   \mathit{acquire}_1(l,z);\ \cstore{m}{1};\ \mathit{release}_1(l) \\
\mathit{consumer} \;=\;\;
 & \cassign{y}{1};\ \opwhile\ y = 1\ \opdo \\
 & \quad (\mathit{acquire}_2(l,y);\ \cload{y}{m}; \\
 & \quad\ \ \opif\ y = 1\ \opthen\
     (\cstore{m}{0};\ \mathit{release}_2(l);\ \cassign{y}{0})\
     \opelse\ (\mathit{release}_2(l);\ \cassign{y}{1})); \\
 & \cload{w}{a_1};\ \dots;\ \cload{w}{a_K}
\end{align*}
The derived specifications are
\begin{align*}
 & \spec{J}{A(\vec{0})}{\mathit{producer}}{\femp} \\
 & \spec{J}{\femp}{\mathit{consumer}}{A(\vec{v})} \\
 & \spec{J}{A(\vec{0})}{\mathit{producer} \oppar \mathit{consumer}}
        {A(\vec{v})}
\end{align*}
(\rocq{mailbox_producer_derived}, \rocq{mailbox_consumer_derived},
\rocq{mailbox_parallel_derived}). The producer starts with the message
cells, all zero, and ends with nothing: it has put them in the mailbox. The
consumer starts with nothing and ends owning the message cells with the
producer's values.

\section{The proof of Section~\ref{sec:daring}}
\label{app:dekker}

The proof uses eleven token cells, none ever read or written by the
program: a prize $w$, and for each thread $k$ the cells $a_k$, $b_k$,
$q_k$, $t_k$ and $p_k$. Below, a token's name stands for the cell
assertion $a_k \fcell 0$, and $\bar{k}$ is the other thread. For each flag
the invariant has three stages,
\begin{align*}
 G_k \;=\;\; & (b_k \fsep q_k \fsep f_k \fcell 0)
     && \text{not yet stored} \\
 \disj\; & (a_k \fsep q_k \fsep (f_k \fcell 0 \fseq f_k \fwrite{k} 1))
     && \text{stored, perhaps still buffered} \\
 \disj\; & (a_k \fsep t_k \fsep f_k \fcell 1)
     && \text{fenced}
\end{align*}
and with $F_k$ for the third of these,
\[
  J \;=\; (w \fsep G_1 \fsep G_2) \;\disj\; (p_1 \fsep F_1 \fsep G_2)
     \;\disj\; (p_2 \fsep F_2 \fsep G_1).
\]
The first disjunct says nobody has taken the prize. The others say thread
$k$ has taken it, leaving $p_k$ in its place, and had fenced first. Thread
$k$ starts with $a_k \fsep t_k \fsep p_k$. Each step is one use of
\rn{atomic}, and the tokens the thread holds determine which disjuncts it
can find:
\begin{center}
\small
\begin{tabular}{@{}lll@{}}
\toprule
Step & Thread $k$ holds afterwards & Exchange with the invariant \\
\midrule
$\cstore{f_k}{1}_k$ & $b_k \fsep t_k \fsep p_k$ &
  gives $a_k$, takes $b_k$; the write joins $G_k$ \\
$\cfence_k$ & $b_k \fsep q_k \fsep p_k$ &
  gives $t_k$, takes $q_k$; the history has drained \\
$\cload{r_k}{f_{\bar{k}}}_k$ &
  $(r_k = 1 \conj b_k \fsep q_k \fsep p_k) \disj
   (r_k = 0 \conj b_k \fsep q_k \fsep w)$ &
  on reading $0$, gives $p_k$ and takes $w$ \\
\bottomrule
\end{tabular}
\end{center}
At the load, a thread still holding $p_k$ cannot find its own winning
disjunct. If the other thread has won, its flag is $1$ and the load reads
$1$. Otherwise the prize is still there, and a thread that reads $0$ takes
it. The second stage of $G_{\bar{k}}$ is where the race is: the load may
find the other thread's write still buffered, and reads $0$ or $1$
accordingly. By the concurrency rule the final state is the spatial
conjunction of the two postconditions, and $w \fsep w$ is false, so
$r_1 = r_2 = 0$ is excluded.

An unlocked memory satisfying $J$ has one of 24 shapes, and each of the six
premises is proved by a checked computation over all of them
(\rocq{buffered_dekker_store0_premise} and its five companions).



\end{document}
